\documentclass[10pt,a4paper]{article}

\usepackage{geometry}
\geometry{a4paper,total={170mm,257mm},left=20mm,top=20mm}
\usepackage[utf8]{inputenc}
\usepackage{amsmath,amsthm,amssymb,mathrsfs,natbib,graphicx}
\usepackage{listings,caption,tikz,array,makecell,sidecap,subcaption}
\usepackage{xcolor,color,enumitem,comment}

\IfFileExists{algorithm2e.sty}{\usepackage[ruled,vlined]{algorithm2e}}{}
\IfFileExists{algpseudocode.sty}{\usepackage{algpseudocode}}{}
\IfFileExists{bbm.sty}{\usepackage{bbm}}{}
\usepackage[most]{tcolorbox}
\usepackage{tabularray,diagbox,csquotes}
\UseTblrLibrary{booktabs}
\setcitestyle{numbers,comma,square}
\definecolor{darkred}{RGB}{139,0,0}
\definecolor{darkblue}{RGB}{0,0,139}
\newif\ifshowchanges
\showchangesfalse
\newcommand{\rev}[1]{{\ifshowchanges\color{darkblue}\fi#1}}
\newcommand{\dr}[1]{\rev{#1}}
\newenvironment{drblock}{\par\begingroup\color{darkblue}}{\endgroup\par}
\usepackage{xurl}
\usepackage[pdfpagelabels,hypertexnames=false,bookmarksopen=true,
 bookmarksopenlevel=1,colorlinks=true,pdfstartview=Fit,
 linkcolor=blue,citecolor=blue,urlcolor=blue]{hyperref}
\hypersetup{pdftitle={Sum-of-Squares and Non-Sum-of-Squares Regimes for the Toeplitz Bottcher-Wenzel Form},
 pdfauthor={Wenqi Zhu and Ping Nie}}
\usepackage{cleveref}

\newtheorem{proposition}{Proposition}[section]
\newtheorem{theorem}{Theorem}[section]
\newtheorem{corollary}{Corollary}[section]
\newtheorem{lemma}{Lemma}[section]
\newtheorem{example}{Example}[section]
\newtheorem{remark}{Remark}[section]
\theoremstyle{definition}
\newtheorem{definition}{Definition}[section]
\newtheorem{assumption}{Assumption}
\renewcommand{\baselinestretch}{1.25}
\setlength{\emergencystretch}{2em}
\allowdisplaybreaks[1]
\newcommand{\R}{\mathbb{R}}
\newcommand{\C}{\mathbb{C}}
\newcommand{\SoS}{\mathrm{SoS}}
\newcommand{\calF}{\mathcal{F}}
\newcommand{\calR}{\mathcal{R}}
\newcommand{\one}{\mathbf{1}}
\newcommand{\norm}[1]{\lVert#1\rVert}
\DeclareMathOperator{\tr}{tr}
\DeclareMathOperator{\diag}{diag}
\DeclareMathOperator{\Span}{span}
\DeclareMathOperator{\rank}{rank}
\newcommand{\ThresholdExponent}{9961475}
\newcommand{\SufficientOrder}{2^{\ThresholdExponent}}
\newcommand{\CornerDepth}{2^{\number\numexpr\ThresholdExponent-1\relax}}
\newcommand{\ScaleExponent}{\number\numexpr(\ThresholdExponent-1)/2\relax}

\title{Sum-of-Squares and Non-Sum-of-Squares Regimes for the Toeplitz Böttcher–Wenzel Form}

\author{Wenqi Zhu\thanks{Mathematical Institute, Woodstock Road, University of Oxford, Oxford, UK, OX2 6GG.  {\tt wenqi.zhu@maths.ox.ac.uk} } \quad and \quad
Ping Nie\thanks{David R. Cheriton School of Computer Science, University of Waterloo.
{\tt ping.nie@uwaterloo.ca}}}

\date{\today}

\begin{document}
\maketitle

\begin{abstract}
Toeplitz matrices are among the simplest structured models of
translation invariance, arising naturally in convolution, stationary
covariance models, and discretizations of translation-invariant
operators.  For a pair of real matrices, the B\"ottcher--Wenzel form
(BW form) is a nonnegative quartic measuring the gap in the
corresponding commutator inequality.  L\'aszl\'o conjectured that,
when both matrices are Toeplitz, this quartic is a sum of squares
for every matrix order. 
We show that the conjecture is false.  More precisely, we identify an
explicit family of corner-supported Toeplitz pairs, retaining only the
outermost diagonals on the two sides, for which the associated BW form
is not a sum of squares at all sufficiently large orders.  The proof
reduces the complete family of Gram representations to a finite complex
test: its Gram-invariant part develops a strictly negative limiting
direction, while the remaining Gram-dependent contribution vanishes.
We also give an explicit sufficient threshold for the matrix order. 
The negative result is complemented by two positive regimes.  If either
Toeplitz factor is symmetric or skew-symmetric, then the corresponding
BW form is SoS for every order, with the other factor allowed to be an
arbitrary real Toeplitz matrix.  For unrestricted Toeplitz pairs, we
further prove SoS representability for every $2\le N\le50$, with
formal Lean~4 verification through order $20$.
\end{abstract}
\section{Introduction}\label{sec:introduction}
For a real vector $x=(x_{1-N},\dots,x_{N-1})$ let $T_N(x)=(x_{i-j})_{0\le i,j<N}$
be the real Toeplitz matrix of order $N$ with symbol $x$. Throughout, norms
and inner products of matrices are Frobenius norms and inner products. Put
\begin{equation}\label{eq:FN}
 F_N(x,y)=2\|T_N(x)\|_F^2\|T_N(y)\|_F^2
 -2\langle T_N(x),T_N(y)\rangle_F^2
 -\|[T_N(x),T_N(y)]\|_F^2,
\end{equation}
a homogeneous quartic in $4N-2$ real variables. The B\"ottcher--Wenzel
inequality implies that $F_N$ is nonnegative. L\'aszl\'o conjectured that
this nonnegativity can always be expressed as a sum of squares of real
quadratic forms \cite[Conjecture~15]{laszlo2012}; the question is problem
MI-15 in the OpenProblemsInNLA catalogue \cite{catalogue2026}. The
following theorem answers it negatively.

\begin{theorem}\label{thm:negative}
For every integer $N\ge\ObstructionOrder$, the polynomial $F_N$ is not a finite
sum of squares of real homogeneous quadratic forms. Nevertheless,
$F_N(x,y)\ge0$ for all real symbol vectors $x,y$ and every $N\ge1$.
\end{theorem}

A \emph{sum of squares} (SOS) always means a finite sum of squares of
real homogeneous quadratic forms, with arbitrary real coefficients and
any number of terms. For a quartic form, requiring the squares to be
homogeneous quadratic forms loses nothing: if a form of degree four
equals $\sum_jp_j^2$ for real polynomials $p_j$, then neither the squares
of the highest-degree homogeneous parts of the $p_j$ nor those of their
lowest-degree parts can cancel, so every $p_j$ is a quadratic form. The
nonnegativity assertion is classical (Section~\ref{sec:context}). The
threshold is sufficient, not optimized; the first order at which a sum of
squares fails is unknown.

\subsection{The obstruction in outline}
Let $\calB_m$ be the restriction of $F_{2m}$ to the $m$ outermost nonzero
diagonals on each side, all other symbol variables being zero. After
relabelling by depth, the same polynomial is a restriction of $F_N$ for
every $N\ge2m$ (Lemma~\ref{lem:corner}). Substitution preserves a sum of
squares, so a single obstruction for $\calB_m$ transfers to every larger
order.

\emph{What is tested.} By Lemma~\ref{lem:frame}, $\calB_m$ is a sum of
squares exactly when $\calB_m=z^TQz$ for a positive semidefinite Gram
matrix $Q$ on the wedge coordinates $z_{ab}=x_ay_b-x_by_a$; such a $Q$
is determined only up to the Pl\"ucker relations among the wedges. The
test is a linear functional $\Phi_m$ of $Q$ that is nonnegative whenever
$Q\succeq0$. To form it, average $Q$ over two symmetries of the corner,
which keeps it positive semidefinite, and split the result into three
positive semidefinite blocks. Substituting monomials in two complex
variables for the wedge coordinates turns each block into a generating
kernel in four complex variables. A fixed combination of these kernels,
evaluated at all pairs of $1024$ points of $\C^2$, is a positive
semidefinite Hermitian $1024\times1024$ matrix, and $\Phi_m$ is its
quadratic pairing with a fixed real vector $c$.

\emph{At which points.} The points and the coefficients are fixed once
and for all:
\begin{equation}\label{eq:witness}
\begin{gathered}
 I=\{(r,j,\sigma):r\in\{1,2\},\ 1\le j\le256,\ \sigma\in\{-1,1\}\},\\
 a_\ell=r+i\sigma j,\quad b_\ell=r-i\sigma j,\quad
 c_\ell=\sigma\hat c_j,\qquad
 \hat c_j=\left\lfloor\frac{2^{43}j}{(128^2+j^2)^2}\right\rfloor.
\end{gathered}
\end{equation}
Thus $|I|=1024$ and $\|c\|_1=859652168<2^{30}$. For a dyadic scale
$\epsilon=2^{-k}$ and a corner of depth $m=\epsilon^{-2}$, the point
attached to $\ell=(r,j,\sigma)$ is
$(ie^{-\epsilon a_\ell},-ie^{-\epsilon b_\ell})$. It weights a diagonal at
depth $p$ by $e^{-\epsilon pr}$ and therefore probes depths of order
$1/\epsilon$, deep inside the corner.

\emph{Why a negative value contradicts SOS.} An exact identity
(\S\ref{sec:mech-identity}) writes the evaluated combination as a
part that is the same for every $Q$ minus a single $Q$-dependent term.
After scaling, the fixed part tends to a limiting tangent kernel whose
pairing with $c$ is below $-2^{42}$ by exact rational arithmetic
(Lemma~\ref{lem:witness}). Positivity of $Q$ forces the $Q$-dependent
term to be small, with constants independent of $Q$, of its size and of
its rank (Sections~\ref{sec:tails} and~\ref{sec:continuation}).
Theorem~\ref{thm:finite-scale} states the precise result: for every
integer $k\ge16$ and $m=2^{2k}$, if $\calB_m$ is a sum of squares, then
\[
 0\le\Phi_m<-2^{42}+2^{90}m^{-1/2}+2^{80}m^{-1/262144}.
\]
The right side is negative once $m$ is large enough. Then no positive
semidefinite Gram matrix, and hence no sum of squares, represents
$\calB_m$.

\begin{corollary}[Explicit threshold]\label{cor:threshold}
The corner $\calB_{m_0}$ is not a sum of squares for
$m_0=\CornerDepth$. Consequently $F_N$ is not a sum of squares for
every $N\ge2m_0=\ObstructionOrder$.
\end{corollary}
\begin{proof}
Set $k=38\cdot2^{17}+1=\ScaleExponent$ in the bound
\eqref{eq:finite-scale} of Theorem~\ref{thm:finite-scale}.
Convexity gives $2^{-x}\le1-x/2$ for $0\le x\le1$, so the right side is
\[
 -2^{42}+2^{90-k}+2^{42}2^{-2^{-17}}
 \le-2^{24}+2^{90-k}<0.
\]
This contradicts $\Phi_{m_0}\ge0$. Substitution preserves a sum of squares,
and the stabilized corner is a restriction of every $F_N$ with
$N\ge2m_0$.
\end{proof}

\subsection{Small orders and the open range}
For every $1\le N\le50$, $F_N$ is a sum of squares
(Proposition~\ref{prop:positive}). Here $F_1=0$; for $2\le N\le50$ the
statement follows from exact rational Gram certificates, and for
$2\le N\le20$ it is also formally proved in Lean~4 with Mathlib. The remaining orders are open:
\[
 51\le N<\ObstructionOrder.
\]
The smallest non-SOS order therefore lies between $51$ and
$\ObstructionOrder$, inclusive; monotonicity of the answer within this
interval is not known.
Section~\ref{sec:discussion} discusses an expectation for this range and the
evidence for it.

\subsection{Nonnegativity and sums of squares}\label{sec:context}
The B\"ottcher--Wenzel inequality
\[
 \|[X,Y]\|_F^2\le2\|X\|_F^2\|Y\|_F^2
\]
holds for arbitrary real square matrices
\cite{bottcher2005,laszlo2007,vong2008,bottcher2008}. If $X\ne0$,
replacing $Y$ by $Y_\perp=Y-\langle X,Y\rangle_F X/\|X\|_F^2$ does not change the
commutator. Expanding $\|Y_\perp\|_F^2$ in the resulting inequality gives
exactly $F_N\ge0$; the case $X=0$ is immediate. Thus the nonnegativity
assertion of Theorem~\ref{thm:negative} holds without the Toeplitz
hypothesis.

A nonnegative quartic need not be a sum of squares of quadratic forms.
This happens already for biquadratic forms, which have degree two in
each of two groups of variables, as $F_N$ has in $x$ and in $y$: Choi
gave a nonnegative biquadratic form in $3+3$ variables that is not a sum
of squares \cite{choi1975}. A form is a sum of squares exactly when it
admits a positive semidefinite Gram matrix, and its Gram matrices form an
affine family \cite{choi1995}; this Gram-matrix method underlies both the
certificates and the obstruction in this paper. L\'aszl\'o formulated
the Toeplitz question after studying the sum-of-squares problem for
the B\"ottcher--Wenzel form of general matrices
\cite{laszlo2010,laszlo2012};
Lu and Wenzel discuss it together with the broader commutator
inequalities \cite{lu2017}. The related DDVV inequalities involve sums
of commutator norms for families of symmetric matrices
\cite{desmet1999,ge2008,lu2011jfa}. Their sum-of-squares questions are
likewise distinct from the inequalities themselves; a recent survey
places both in this setting \cite{ge2024published}.

Theorem~\ref{thm:negative} shows that the Toeplitz structure does not
close the gap between the two properties. A sum of squares requires a
positive semidefinite Gram matrix on the whole space of wedge
coordinates, whereas nonnegativity only tests decomposable wedges
$x\wedge y$. The proof finds a negative direction for the stronger
requirement in a limiting kernel and then controls every error at a
finite order.

\subsection{Verification}
The negative conclusion of Theorem~\ref{thm:negative} is formally proved
in Lean~4 with Mathlib, for the literal polynomial $F_N$ and arbitrary
finite sums of real homogeneous quadratic squares, with no remaining
analytic hypothesis. The proof includes the complete corner extraction,
the analytic estimates, and a kernel-checked exact witness
(Lemma~\ref{lem:witness}). Its expanded statement is
\path|ToeplitzSOS.Negative.sharp_negative_explicit|; the negation of the
catalogue conjecture is \path|ToeplitzSOS.Negative.not_MI15|.
The positive range is formally proved through order $20$.
Appendix~\ref{app:verification} gives the declarations, the scope of each
formal component, and the relation to the written proof below.

\section{Failure of a sum-of-squares representation for Toeplitz BW Form}\label{sec:non-sos}

In this section, we prove the following theorem.
\begin{tcolorbox}[colback=white,colframe=black!60,boxrule=0.5pt,
 arc=0pt,left=7pt,right=7pt,top=3pt,bottom=3pt]
\begin{theorem}\label{thm:negative}
There exists an integer $N_0$ such that, for every integer $N\ge N_0$,
the polynomial $F_N$ in \eqref{eq:FN} is nonnegative but is not a finite
sum of squares of real homogeneous quadratic polynomials.
An explicit sufficient value of $N_0$ is given in
\Cref{app:explicit-bound}.
\end{theorem}
\end{tcolorbox}

This section is organized as follows.  In \Cref{subsec:wedge-coordinates},
following \cite[Section~2]{laszlo2012}, we introduce the wedge coordinates
and the associated wedge Gram representation.  These coordinates will be
used throughout the paper.  In \Cref{subsec:corner}, we introduce a class
of corner-supported Toeplitz matrices.  The B\"ottcher--Wenzel forms
generated by this class will provide the non-SoS examples at sufficiently
large matrix size.

Our proof proceeds in three steps.  We first define the corresponding
corner-supported B\"ottcher--Wenzel form, denoted by $\mathcal F_m$, and
construct a convenient Gram representation for it (\Cref{subsec:corner-gram}).  Since Gram representations are not unique, we
then use the symmetries of $\mathcal F_m$ to reduce the full Gram family
to a general normal form (\Cref{lem:normal-form}).  Finally, we construct a finite matrix test which, for
sufficiently large $m$, rules out positive semidefiniteness for every
possible Gram representation of $\mathcal F_m$ (\Cref{subsec:finite-test}).  This proves that
$\mathcal F_m$ is not SoS for sufficiently large $m$, and hence, that $F_N$ is not SoS for all sufficiently large $N$.

\subsection{Wedge-coordinate and wedge Gram representation}\label{app:wedge-gram}\label{subsec:wedge-coordinates}

Following \cite[Section~2]{laszlo2012}, we use the wedge-coordinate representation of \eqref{eq:FN} throughout the paper. The BW form of Toeplitz matrices $X,Y$ is a nonnegative homogeneous quartic polynomial in the $4N-2$ real variables
\begin{equation}
    \label{origin parameter}
(x_{-(N-1)},\ldots,x_{N-1},y_{-(N-1)},\ldots,y_{N-1}) \in\R^{4N-2} 
\end{equation}
with $x,y\in\R^{2N-1}$. 
We define the wedge coordinates of $(x,y)$ by
\begin{equation}
z_{ab}=x_ay_b-x_by_a,\qquad -(N-1)\le a<b\le N-1,
    \label{wedge coordinate}
\end{equation}
where $x_a$ denotes the entry of $x$ indexed by $a$ and $y_b$ denotes the entry of $y$ indexed by $b$. Collect these coordinates, in \begin{equation}
z:=\bigl(z_{-(N-1),-(N-2)},z_{-(N-1),-(N-3)},\ldots,z_{-(N-1),N-1},
z_{-(N-2),-(N-3)},\ldots,z_{N-2,N-1}\bigr)^\top
\in\R^{\binom{2N-1}{2}}.
\end{equation}
The vector $z$ is a quadratic lifting of the original variables.  
The wedge coordinates should be viewed as a structured quadratic lifting of the original Toeplitz parameters in \eqref{origin parameter} rather than as a linear change of variables.  Each Toeplitz matrix is determined by $2N-1$ real parameters, so the pair $(X,Y)$ is described by $4N-2$ real variables.  The wedge vector
$
z=(z_{ab})_{a<b}, 
z_{ab}=x_ay_b-x_by_a,
$
has
$
\binom{2N-1}{2}
$
coordinates, each of which is quadratic in the original variables.  Thus the map
$
(x,y)\mapsto z=x\wedge y
$
lifts the quartic form $F_N(x,y)$ to a quadratic form in the wedge coordinates,
$
F_N(x,y)=z^\top Qz.
$
The wedge representation is substantially smaller than the full degree-two monomial representation in the $4N-2$ original variables, whose dimension is
$
\binom{4N-1}{2}.
$
The wedge coordinates are not algebraically independent; their quadratic relations are precisely the Pl\"ucker relations described in \Cref{lem:frame}.  Thus the wedge coordinates form a structured subcollection of the usual degree-two monomial representation rather than an invertible change of coordinates.



We call the real symmetric matrix $Q\in\mathbb S^{\binom{2N-1}{2}}$ a wedge Gram matrix of $F_N$ if
$$
        F_N(x,y)=z^\top Qz.
$$
% The polynomial variables and Gram matrices remain real.\footnote{Complex test vectors are used later: a real symmetric positive semidefinite $Q$ also satisfies $v^*Qv\ge0$ for $v\in\C^d$. This does not change the real polynomial identity $F_N=z^\top Qz$ into a sesquilinear identity in complex symbol variables.}
Such a Gram matrix is generally not unique; \Cref{lem:frame} characterizes this full Gram family.
It is well known that a polynomial is SoS if and only if it admits a positive semidefinite Gram representation \cite{parrilo2000structured,lasserre2001global}. For the B\"ottcher--Wenzel form, \Cref{lem:frame} also shows that every quadratic factor in an SoS representation lies in the wedge span. Consequently,
$$
        F_N(x,y)\ \text{is SoS}
        \quad\Longleftrightarrow\quad
        F_N(x,y)=z^\top Qz
        \ \text{for some }Q\succeq0.
$$
% More precisely, if $\tilde{v}_{xy}$ denotes the vector of bilinear monomials $x_a y_b$, then there is a fixed signed matrix $S_N$ such that $z=S_N\tilde{v}_{xy}$. Let
% $\tilde{v}_{xy}$ denotes the vector of bilinear monomials $x_a y_b$ such that 

% Note that \Cref{lem:frame} gives a general formulation
% of this reduction and characterizes all admissible Gram
% representations. The lemma shows that these are not merely some Gram
% redundancies: they describe \emph{all} Gram redundancies relevant to the
% quartic forms considered below.


\begin{lemma}[Complete characterization of wedge Gram representations]\label{lem:frame}
Let $I$ be a finite ordered set. Let \(F(x,y)\) be a real homogeneous quartic satisfying
\begin{equation}
    \label{vanish}
    F(x,0)=F(0,y)=F(x,x)=0
    \qquad
    \text{for all }x,y\in\R^I.\end{equation}
 If $F$ is a sum of squares, every quadratic factor is
a linear combination of the $z_{ab}$.
Moreover, the real symmetric corrections $\mathcal E$ satisfying
$z^\top\mathcal E z\equiv0$ are spanned by the matrices
$P_{abcd}$ specified by\footnote{Precisely, its entries at $(ab,cd),(ac,bd),(ad,bc)$ are $1,-1,1$, with the same entries at transposed positions and zero elsewhere.}
\begin{equation}\label{eq:plucker}
 z^\top P_{abcd}z
 =2(z_{ab}z_{cd}-z_{ac}z_{bd}+z_{ad}z_{bc}),
 \qquad a<b<c<d.
\end{equation}
Consequently, if $F=z^\top Gz$, then
\begin{equation}\label{eq:frame}
 F\text{ is SoS}
 \quad\Longleftrightarrow\quad
 G+\mathcal E\succeq0\text{ for some }\mathcal E\in\Span\{P_{abcd}\}.
\end{equation}
\end{lemma}
\begin{proof}

Write $F=\sum_jq_j^2$. Evaluating at $y=0$ shows that every
$q_j(x,0)$ vanishes identically, and evaluating at $x=0$ gives the
analogous conclusion. Thus $q_j=x^\top A_jy$ for a real matrix $A_j$.
The identity $q_j(x,x)=0$ implies $A_j+A_j^\top=0$, whence
$q_j=\sum_{a<b}(A_j)_{ab}z_{ab}$. The coefficient vectors of the
$q_j$ give a positive semidefinite Gram matrix. 
To determine the relation space, group products of wedges by their
multisets of labels. Monomials belonging to different multisets cannot
cancel. For a multiset with a repeated label, there is at most one
nonzero wedge product up to sign, and that product is not the zero
polynomial. For four distinct labels, the three possible products are
$z_{ab}z_{cd}$, $z_{ac}z_{bd}$ and $z_{ad}z_{bc}$. Any two are linearly
independent, as is seen by comparing their coefficients, whereas their
alternating combination in \eqref{eq:plucker} vanishes. Their relation
space is therefore exactly one-dimensional. This proves the asserted
description of the correction space. 
Any representing Gram matrix differs from $G$ by an element of this
space. Conversely, a positive semidefinite representative factors as
$R^\top R$, and the rows of $Rz$ provide the required quadratic factors.
This proves \eqref{eq:frame}.
\end{proof}

Note that Lemma~\ref{lem:frame} does not require Toeplitz structure. It applied to the Toeplitz Bottcher--Wenzel form since for the polynomial $F_N$, the three vanishing conditions
\eqref{vanish} hold.

\subsection{Toeplitz subclass that is non-SoS for large $N$}
\label{subsec:corner}

We begin by identifying a Toeplitz subclass on which the large-order
failure of SoS representability will be detected.  Fix a positive integer $m$ and
 $N\ge 2m$.  We call a Toeplitz matrix
\emph{corner-supported at depth $m$} if only the outermost $m$
diagonals on each side are retained and all Toeplitz coefficients in
the central band are set to zero.
The following small example illustrates the geometry.
Take $N=6$ and $m=2$, a corner-supported Toeplitz matrix with depth $2$ has the form
\begin{equation}\label{eq:corner-example-matrix}
X_{(N=6,m=2)}
=
\begin{pmatrix}
0&0&0&0&x_{-4}&x_{-5}\\
0&0&0&0&0&x_{-4}\\
0&0&0&0&0&0\\
0&0&0&0&0&0\\
x_4&0&0&0&0&0\\
x_5&x_4&0&0&0&0
\end{pmatrix}.
\end{equation}
Note that this example is only to illustrate the corner
geometry.  At this small depth $m=2$, the corresponding BW form is
still SoS; the failure of SoS representability occurs only for sufficiently large corner
depth. 

More generally, let $X$ and $Y$ be corner-supported Toeplitz matrices with at depth $m$, then we denote the indices of the retained entry (non-zero band) of Toeplitz labels as
\[
 \Lambda_{N,m}
 :=
 \{\pm(N-1), \pm (N-2), \dotsc, \pm(N-m)\}
\]
such that $ x_a=y_a=0$ for $a \notin \Lambda_{N,m}$.
The \eqref{eq:FN} with respect to in corner-supported $X,Y$ is the restriction to the the quartic
\begin{equation}\label{eq:Bm-definition}
 \mathcal F_m
 :=
 F_N(x,y)\Big|_{
 x_a=y_a=0\;(a\notin\Lambda_{N,m})}.
\end{equation}

\begin{lemma}[Properties of the corner-supported Toeplitz BW form]\label{lem:corner}
Let $X$ and $Y$ be corner-supported Toeplitz matrices of size $N$ and
depth $m$, with $N\ge 2m$, and let $\mathcal F_m$ be the restricted
B\"ottcher--Wenzel form defined in \eqref{eq:Bm-definition}.  Let
$
 z=(z_{ab})_{a<b},
 z_{ab}:=x_a y_b-x_b y_a,
$
be the full wedge vector defined in \eqref{wedge coordinate}, and let
\begin{equation}
\label{restricted wedge coordinate}
 z_r:=(z_{ab})_{
 a,b\in\Lambda_{N,m},\,a<b}
 \in\R^{\binom{2m}{2}}
\end{equation}
be its restriction to the retained labels.
Then the following hold.
\begin{enumerate}
\item
If
$
        F_N=z^\top Gz
$
is any Gram representation in the full wedge coordinates, then
\[
        \mathcal F_m=z_r^\top G_rz_r,
\]
where $G_r$ is the principal submatrix of $G$ indexed by the retained
wedge coordinates.

\item
After identifying the retained labels by their side and depth
$p=0,\ldots,m-1$, the polynomial $\mathcal F_m$ depends only on $m$ and
is independent of $N$.

\item
If $\mathcal F_m$ is not SoS, then $F_N$ is not SoS for every
$N\ge2m$.
\end{enumerate}
\end{lemma}

\begin{proof}
For the first statement, setting $x_a=y_a=0$ for $a\notin\Lambda_{N,m}$ annihilates every wedge containing a non-retained label. Hence $z^\top Gz$ restricts directly to $z_r^\top G_rz_r$.

For the second statement, let $U_a$ be the shift matrix $(U_a)_{ij}=\one_{i-j=a}$, $0\le i,j<N$. Distinct shifts are Frobenius orthogonal and $\|U_a\|_F^2=N-|a|$. The weighted Lagrange identity \cite[Lemma~10]{laszlo2012} and bilinearity of the commutator give
\begin{equation}\label{eq:literal-gram}
 F_N=
 2\sum_{a<b}(N-|a|)(N-|b|)z_{ab}^2
 -
 \Bigl\|\sum_{a<b}z_{ab}[U_a,U_b]\Bigr\|_F^2.
\end{equation}
For a retained label, $N-|\pm(N-1-a)|=a+1$. Shifts on the same side commute, while for $C_{ab}:=[U_{-(N-1-a)},U_{N-1-b}]$ one has, for $N\ge2m$,
\begin{equation}\label{eq:commutator-overlap}
 \langle C_{ab},C_{cd}\rangle_F
 =
 2\one_{a-b=c-d}\min(a+1,b+1,c+1,d+1).
\end{equation}
Thus every coefficient of the restricted quartic depends only on the depths $a,b,c,d$, and hence only on $m$.

Finally, linear restriction preserves sums of squares. Therefore an SoS representation of $F_N$ would restrict to an SoS representation of $\mathcal F_m$, proving the last statement by contraposition.
\end{proof}

\subsubsection{Gram representation of the corner-supported Toeplitz BW form}
\label{subsec:corner-gram}

In this subsection, we study the Gram representations of the
corner-supported Toeplitz B\"ottcher--Wenzel form $\mathcal B$.  We
first construct one fixed Gram representation with a particularly simple
block structure, and then derive a general normal form that captures all
positive semidefinite Gram representations under the reduced restricted
wedge coordinates.

To obtain the structured Gram representation, we first rearrange the
restricted wedge vector as
$
z_r=(u^\top,v^\top,w^\top)^\top,
$
where
$
u,v\in\R^{\binom m2}
$
and
$
w\in\R^{m^2}.
$
The entries are defined by
\[
u_{ab}:=x_ay_b-x_by_a,\quad a<b<0,\qquad
v_{ab}:=x_ay_b-x_by_a,\quad 0<a<b, \qquad
w_{ab}:=x_ay_b-x_by_a,\quad a<0<b.
\]
We refer to $u$ and $v$ as the two \emph{pure wedge} families, since the
two labels in each wedge lie on the same side of the main diagonal, and
to $w$ as the \emph{mixed wedge} family, since its two labels lie on
opposite sides. With this ordering, the restricted Gram representation takes a highly
regular block form.  The structure is already visible in the case
$m=2$.



\begin{example}[Gram representation for a corner-supported Toeplitz BW form with $m=2$]\label{ex:gram-m2}
In~\eqref{eq:corner-example-matrix}, with the ordering described above, the corresponding restricted wedge vector and baseline Gram matrix are
\[
z_r=
\left(
\begin{array}{c}
u_{-5,-4}\\
v_{4,5}\\
w_{-5,5}\\
w_{-5,4}\\
w_{-4,5}\\
w_{-4,4}
\end{array}
\right)
\qquad
\begin{array}{c}
\left.\vphantom{\begin{array}{c}u_{-5,-4}\end{array}}\right\}\ \text{$-$ pure}\\[-0.2ex]
\left.\vphantom{\begin{array}{c}v_{4,5}\end{array}}\right\}\ \text{$+$ pure}\\[-0.2ex]
\left.\vphantom{\begin{array}{c}w_{-5,5}\\ w_{-5,4}\\ w_{-4,5}\\ w_{-4,4}\end{array}}\right\}\ \text{mixed} \\[-0.2ex]
\end{array}
\qquad
G_2=
\left(
\begin{array}{c|c|cccc}
4&2&0&0&0&0\\ \hline
2&4&0&0&0&0\\ \hline
0&0&0&0&0&0\\
0&0&0&2&-2&0\\
0&0&0&-2&2&0\\
0&0&0&0&0&4
\end{array}
\right).
\]
Here the mixed coordinates are ordered by corner depth,
$
(w_{00},w_{01},w_{10},w_{11})
=
(w_{-5,5},w_{-5,4},w_{-4,5},w_{-4,4}).
$
Thus,
\[
\mathcal B_2=z_r^\top G_2z_r
=
2(u_{-5,-4}+v_{4,5})^2
+2u_{-5,-4}^2
+2v_{4,5}^2
+2(w_{-5,4}-w_{-4,5})^2
+4w_{-4,4}^2.
\]
\end{example}
Example~\ref{ex:gram-m2} shows that this ordering arranges the restricted Gram matrix into a particularly simple block structure. The next lemma proves that the same structure holds for every corner depth $m$. 
For notational simplicity, in the rest of the paper,  we index
these three families 
$
u_{ab}:=u_{-(N-1-a),-(N-1-b)},
$
$
v_{ab}:=v_{N-1-b,N-1-a},
$
and
$
w_{ab}:=w_{-(N-1-a),N-1-b}.
$
This is only a relabelling of the same restricted wedge coordinates.


\begin{lemma}[Gram representation of the corner-supported Toeplitz BW form]\label{lem:corner-gram}
Let $\calF_m$ be a corner-supported Toeplitz B\"ottcher--Wenzel form
generated by two Toeplitz matrices $X,Y$ of size $N$ and corner depth
$m$, where $N\ge2m$.  
Collect the coordinates into
$z_r=(u^\top,v^\top,w^\top)^\top$, using lexicographic
order within each depth-indexed group.
% \footnote{The example uses the original-label ordering; the depth ordering here differs by a coordinate permutation. The corresponding Gram matrices are permuted in the same way.} 
Then
\begin{equation}\label{eq:baseline-gram}
 \mathcal F_m=z_r^\top G_mz_r,\qquad
 G_m=
 \begin{pmatrix}
 D&K&0\\
 K&D&0\\
 0&0&B
 \end{pmatrix}.
\end{equation}
Here $D,K\in\mathbb S^{\binom m2}$,
$B\in\mathbb S^{m^2}$, and
$G_m\in\mathbb S^{m(2m-1)}$ where $
\mathbb S 
$
denotes the space of real symmetric matrices and the specific entries are
\begin{equation}\label{eq:baselines}
\begin{aligned}
 (D)_{ab,cd}
 &=2(a+1)(b+1)\delta_{ac}\delta_{bd}, \quad
 (K)_{ab,cd}
  =2\one_{b-a=d-c}\min(a+1,c+1),\\
 (B)_{ab,cd}
 &=2(a+1)(b+1)\delta_{ac}\delta_{bd}
 -2\one_{a+b=c+d}\min(a+1,b+1,c+1,d+1).
\end{aligned}
\end{equation}
For $D,K$, the index pairs satisfy $0\le a<b<m$ and
$0\le c<d<m$; for $B$, all four indices range independently over
$\{0,\ldots,m-1\}$.  In particular, $D,K,B$, and hence $G_m$, are fixed by \eqref{eq:baselines} and are independent of $N$.
% \footnote{The baseline $G_m$ is one fixed choice; it is not uniquely determined by the polynomial. The admissible Gram corrections remain nonunique.}
\end{lemma}

\begin{proof}
We rearrange the restricted expression \eqref{eq:literal-gram} using
the Pl\"ucker relations, as in the Gram construction of
\cite[Sections~2 and~5]{laszlo2012}.  The raw Gram couples mixed wedges
with equal differences $a-b=c-d$.  For $a<c$ and $b<d$ satisfying this
equality, the corresponding term is
$-4\min(a+1,b+1)w_{ab}w_{cd}$.  With the orientation above, the Pl\"ucker
identity becomes
\begin{equation}\label{eq:exchange}
 u_{ac}v_{bd}+w_{ab}w_{cd}-w_{ad}w_{cb}=0,
\end{equation}
so that the mixed product is replaced by
$-u_{ac}v_{bd}+w_{ad}w_{cb}$.  Summing the resulting contributions gives
\eqref{eq:baselines}.
\end{proof}

The Gram representation in Lemma~\ref{lem:corner-gram} is not unique.  By \Cref{lem:frame}, every other representation is obtained by adding Pl\"ucker corrections.  The full Gram family is large, but the symmetries of $\mathcal F_m$
allow us to average any positive semidefinite representative into a
simple normal form.  
% The realignment map $\calR$ appearing below is the index re-pairing defined in the proof.

\begin{lemma}[General Gram normal form for the corner-supported Toeplitz BW form]
\label{lem:normal-form}
Let $\calF_m$ be a corner-supported Toeplitz B\"ottcher--Wenzel form
generated by two Toeplitz matrices $X,Y$ of size $N$ and corner depth
$m$, where $N\ge2m$.  Then the following statements hold.
\begin{itemize}

\item The Gram representation of $\calF_m$ is not unique.  If $\calF_m$
admits a positive semidefinite Gram representation, then it admits a
symmetry-reduced representative of the form
\begin{equation}\label{eq:general-gram-1}
 Q_{\rm av}=G_m+
 \begin{pmatrix}
 \mathcal E_0&\mathcal E_1&0\\
 \mathcal E_1&\mathcal E_0&0\\
 0&0&\calR\mathcal E_1
 \end{pmatrix}
 =
 \begin{pmatrix}
 D+\mathcal E_0&K+\mathcal E_1&0\\
 K+\mathcal E_1&D+\mathcal E_0&0\\
 0&0&B+\calR\mathcal E_1
 \end{pmatrix}\succeq0,
\end{equation}
where $\mathcal E_0,\mathcal E_1\in\mathbb S^{\binom m2}$,
$\mathcal E_0$ is a same-side Pl\"ucker correction acting within each
pure wedge family (equivalently, a pure four-form),
$\mathcal E_1$ is a real symmetric matrix coupling the two pure wedge
families, and $\calR$ is the realignment map induced by the corresponding
index re-pairing.\footnote{The full correction matrix is an element
$\mathcal E$ of the space characterized in Lemma~\ref{lem:frame};
$\mathcal E_0$ and $\mathcal E_1$ are the components remaining after
the symmetry reduction. The realignment operator
$
\calR:\mathbb S^{\binom m2}\to\mathbb S^{m^2}
$
is a linear map that re-pairs the four indices of a pure-wedge Gram block into the corresponding mixed-wedge indices.  In particular, it maps a correction on the pure wedge space to the induced correction on the mixed wedge space. The induced mixed correction
$\calR\mathcal E_1$ belongs to $\mathbb S^{m^2}$.  See
\Cref{app:normal-form-proof} for the detailed construction.}

\item Under the orthogonal change of coordinates
$(u,v,w)\mapsto((u+v)/\sqrt2,(u-v)/\sqrt2,w)$, the Gram representation
in \eqref{eq:general-gram-1} becomes block diagonal:
\begin{equation}\label{eq:general-gram-2}
 \widehat Q_{\rm av}
 :=U^\top Q_{\rm av}U
 =
 \begin{pmatrix}
 P_+&0&0\\
 0&P_-&0\\
 0&0&M
 \end{pmatrix},
 \tag{Block-diag Gram}
\end{equation}
where $U$ is the orthogonal matrix representing the change of coordinates 
% with $r=\binom m2$,
% \[
%  U=
%  \begin{pmatrix}
%  \frac1{\sqrt2}I_r&\frac1{\sqrt2}I_r&0\\
%  \frac1{\sqrt2}I_r&-\frac1{\sqrt2}I_r&0\\
%  0&0&I_{m^2}
%  \end{pmatrix},
%  \qquad U^\top U=I,
% \]
and
\begin{equation}\label{eq:sectors}
 P_\pm=D+\mathcal E_0\pm K\pm\mathcal E_1
 \in\mathbb S^{\binom m2},
 \qquad
 M=B+\calR\mathcal E_1\in\mathbb S^{m^2}.
\end{equation}

\item The corner-supported form satisfies
\[
 \calF_m\text{ is SoS}
 \quad\Longleftrightarrow\quad
 \text{there exist }\mathcal E_0,\mathcal E_1
 \text{ such that }Q_{\rm av}\succeq0
 \quad\Longleftrightarrow\quad
 P_+\succeq0,\quad P_-\succeq0,\quad M\succeq0.
\]
\end{itemize}
\end{lemma}

\begin{proof}
The three statements follow from the complete Pl\"ucker description of
the Gram freedom, the symmetries of $\calF_m$, and the orthogonal
decomposition of the two pure wedge families.  A detailed proof is
given in \Cref{app:normal-form-proof}.
\end{proof}

\subsubsection{A finite Gram test}\label{subsec:finite-test}

We now reduce the search over the nonunique Gram family to a finite-dimensional
matrix test; we denote the matrix as $ \mathcal G_m$.  The purpose of the construction is to separate the test matrix
into a part depending only on the fixed polynomial data and a remaining part
depending on the choice of Gram representation (i.e., $  \mathcal G_m=\mathbf L_m-\mathbf R_m $).  We will show that the first
part has a negative limiting direction, whereas the second part is too small
to compensate for it. 

\begin{lemma}[Finite Gram test]\label{lem:finite-test}
Let $V\in\C^{m(2m-1)\times s}$ be an arbitrary complex matrix and let
$\mathcal G_m$ be a complex Hermitian\footnote{$\mathbb H^s:=\{A\in\mathbb C^{s\times s}:A^*=A\}$ represents the space of complex Hermitian $s\times s$ matrices.} $s\times s$ matrix defined as
\begin{equation}\label{eq:finite-compression}
 \mathcal G_m:=V^*\widehat Q_{\rm av}V\in\mathbb H^s.
\end{equation}
If $\mathcal F_m$ is SoS, then $\mathcal G_m\succeq0$ for every $V$.  In particular,
$c^*\mathcal G_mc\ge0$ for every $c\in\C^s$.
\end{lemma}

\begin{proof}
If $\mathcal F_m$ is SoS, Lemma~\ref{lem:normal-form} gives a normal-form
Gram representation with $\widehat Q_{\rm av}\succeq0$.  Hence, for every
$\hat y\in\C^s$,
\[
 \hat y^*\mathcal G_m\hat y
 =\hat y^*V^*\widehat Q_{\rm av}V\hat y
 =(V\hat y)^*\widehat Q_{\rm av}(V\hat y)\ge0.
\]
Thus $\mathcal G_m\succeq0$, and the last assertion follows by taking
$\hat y=c$.
\end{proof}

\begin{remark}[Complex test points]
\label{rem:complex-test}
The original SoS problem and all Gram representations above remain over
the real numbers.  Complex numbers are introduced only in the test
points used below, and do not change the underlying real SoS problem.

Indeed, let $Q$ be a real symmetric positive semidefinite matrix.  For
any complex vector $z=\hat u+i\hat v$, with
$\hat u,\hat v$ real, we have
\[
 z^*Qz
 =
 (\hat u^\top-i\hat v^\top)Q(\hat u+i\hat v)
 =
 \hat u^\top Q\hat u+\hat v^\top Q\hat v\ge0,
\]
where the cross terms cancel because $Q$ is symmetric.  Thus a real
positive semidefinite Gram matrix remains positive semidefinite when
tested against complex vectors.

We may therefore choose several complex test points and arrange their
corresponding feature vectors as the columns of a complex test matrix
$\mathcal V$.  The resulting matrix
$
\mathcal V^*Q\mathcal V
$
must still be positive semidefinite whenever $Q\succeq0$.  Allowing
complex test points gives us a larger and more flexible space of test
directions, without enlarging the class of admissible Gram
representations. 
In the construction in \Cref{sec choose direction}, we  will choose complex points that approach
$(i,-i)$ along prescribed directions as $m$ increases.  The associated
test matrix develops a strictly negative limiting direction $c$.  Since any
real SoS Gram representation would remain positive semidefinite under
the same complex test, this eventually yields the contradiction used to
prove that the corner-supported form is not SoS.
\end{remark}

\begin{lemma}[Finite Gram decomposition]\label{lem:finite-decomposition}
Assume $V$ contains a special block structure, namely
\begin{equation}\label{eq:special-V}
V=
\left(
\begin{array}{ccc}
V_1^*/\sqrt2&
V_1^*/\sqrt2&
V_3^*
\end{array}
\right)^*
\in\C^{m(2m-1)\times s},
\end{equation}
where $V_1 \in \C^{\binom m2 \times s}$ and
$V_3 \in \C^{m^2\times s}$.
Then, the finite Gram matrix in \eqref{eq:finite-compression} admits the decomposition
\begin{eqnarray}\label{eq:finite-decomposition}
   \mathcal G_m &=& \mathbf L_m-\mathbf R_m
   \\ &=&
\underbrace{
 V_1^*DV_1
 +
 V_3^*
 \bigl(B+\calR(D-K)\bigr)V_3
}_{:= \mathbf L_m \text{ Gram-invariant}}
-
\underbrace{
 V_3^*(\calR P_-)V_3
}_{:= \mathbf R_m \text{ Gram-dependent}} .
\end{eqnarray}
Here, \emph{Gram-invariant} means that $\mathbf L_m$ depends only on the fixed polynomial data $D,K,B$ and is independent of the particular Gram representation, whereas \emph{Gram-dependent} means that $\mathbf R_m$ depends on the chosen Gram representative through $P_-$.
\end{lemma}

\begin{proof}
From the block diagonal form of $\widehat Q_{\rm av}$,
\[
\begin{aligned}
\mathcal G_m
=
\left(
\begin{array}{ccc}
V_1^*/\sqrt2&
V_1^*/\sqrt2&
V_3^*
\end{array}
\right)
\left(
\begin{array}{ccc}
P_+&0&0\\
0&P_-&0\\
0&0&M
\end{array}
\right)
\left(
\begin{array}{c}
V_1/\sqrt2\\
V_1/\sqrt2\\
V_3
\end{array}
\right)=
V_1^*\frac{P_++P_-}{2}V_1
+
V_3^*MV_3.
\end{aligned}
\]
Using \eqref{eq:sectors},
$
(P_++P_-)/2=D+\mathcal E_0
$
and
$
M=B+\calR\mathcal E_1.
$
Moreover,
$
P_-=D-K+\mathcal E_0-\mathcal E_1
$
 and $\calR \mathcal E_0=-\mathcal E_0$\footnote{By definition, realignment exchanges the middle two indices,
$
 (\calR \mathcal E_0)_{ac,bd}
 =(\mathcal E_0)_{ab,cd}
 =\Omega_{abcd}.
$
Since $\Omega$ is totally antisymmetric,
$
 \Omega_{acbd}=-\Omega_{abcd}.
$
Equivalently,
$
 (\calR \mathcal E_0)_{ac,bd}
 =-(\mathcal E_0)_{ac,bd}.
$
Hence $\calR \mathcal E_0=-\mathcal E_0$.}.
% \footnote{This is an equality of scalar generating expressions, not of pure- and mixed-space matrices, which have different sizes. Realignment exchanges the middle two arguments of the totally alternating generating expression, changing its sign.}. 
At the level of the corresponding finite evaluations,
this gives
$
V_1^*\mathcal E_0V_1=-V_3^*(\calR \mathcal E_0)V_3.
$
Hence
\[
\begin{aligned}
\mathcal G_m
&=
V_1^*DV_1
+
V_3^*BV_3
-
V_3^*\calR(\mathcal E_0-\mathcal E_1)V_3\\
&=
V_1^*DV_1
+
V_3^*
\bigl(B+\calR(D-K)\bigr)V_3
-
V_3^*(\calR P_-)V_3,
\end{aligned}
\]
which proves \eqref{eq:finite-decomposition}.
\end{proof}

The importance of Lemma~\ref{lem:finite-decomposition} is that
$\mathbf L_m$ depends only on the fixed matrices $D,K,B$, and hence is
independent of the chosen Gram representation.  The entire remaining
dependence on the nonunique Gram representation is contained in
$\mathbf R_m$.  We will choose a finite family of points for which
$\mathbf L_m$ develops a negative direction, while the contribution of
$\mathbf R_m$ vanishes after the natural scaling.

\subsubsection{Choice of the finite test points}
\label{sec choose direction}

We now choose a specific family of test matrices $V_m$ and a direction
$c$ for which the finite compression
$
\mathcal G_m=V_m^*\widehat Q_{\rm av}V_m
$
cannot remain positive semidefinite when $m$ is sufficiently large.
% The construction is designed so that, in the decomposition
% $
% \mathcal G_m=\mathbf L_m-\mathbf R_m,
% $
% the Gram-invariant part $\mathbf L_m$ develops a strictly negative
% direction, while the Gram-dependent part $\mathbf R_m$ becomes
% asymptotically negligible in the same direction.  Since every SoS
% representation would require $\mathcal G_m\succeq0$, such a choice of
% $V_m$ and $c$ rules out all positive semidefinite Gram representations
% of $\mathcal F_m$.

Following the definition of $V$ in \eqref{eq:special-V}. We next choose $V$ so that the three diagonal blocks of
$\widehat Q_{\rm av}$ are tested on the following
coordinates.    For points
$\mathsf x_\ell=(\xi_\ell,\eta_\ell)\in\C^2$, $1\le\ell\le s$, define
\begin{equation}\label{eq:matched-features}
V_1:=
\left[
 \bigl(\xi_\ell^a\eta_\ell^b-\xi_\ell^b\eta_\ell^a\bigr)_{0\le a<b<m}
\right]_{\ell=1}^s
\in\C^{\binom m2 \times s},
\qquad
V_3:=
\left[
 \bigl(\xi_\ell^a\eta_\ell^b\bigr)_{0\le a,b<m}
\right]_{\ell=1}^s
\in\C^{m^2\times s},
\end{equation}
For each $m$, define
\begin{equation}\label{eq:test-nodes}
 \mathsf x_\ell^{(m)}
 =
 \bigl(\xi_\ell^{(m)},\eta_\ell^{(m)}\bigr)
 :=
 \left(
 ie^{-a_\ell/\sqrt m},
 -ie^{-\bar a_\ell/\sqrt m}
 \right),
 \qquad 1\le\ell\le s.
\end{equation}
where the parameters begin with
\[
\begin{alignedat}{6}
a_1&=1+i,      &\qquad a_2&=1-i,      &\qquad a_3&=2+i,      &\qquad a_4&=2-i,      &\qquad a_5&=1+2i,      &\qquad a_6&=1-2i,\\
a_7&=2+2i,     &\qquad a_8&=2-2i,     &\qquad a_9&=1+3i,     &\qquad a_{10}&=1-3i,  &\qquad a_{11}&=2+3i,  &\qquad a_{12}&=2-3i,\ \dotsc
\end{alignedat}
\]
and the pattern continues in conjugate pairs on the two vertical lines
$\Re a=1$ and $\Re a=2$. 

Note that $\mathsf x_\ell^{(m)}\to(i,-i)$ as $m\to\infty$.  From now on
$V_1,V_3,V,\mathbf L_m$ and $\mathbf R_m$ refer to the points
in \eqref{eq:test-nodes}.
The choice of $\Re a=1,2$ is not the only negative example.  Additional
families may also be included in the finite test (see
Appendix~\ref{C:rem:alternative-test-families}).

\begin{proposition}[Negative limit of the Gram-invariant part]
\label{lem:negative-limit}
Let $\mathbf L_m$ and $\mathbf R_m$ be defined as in
\Cref{lem:finite-decomposition}.  Fix the $s$ test points in
\eqref{eq:test-nodes}.  Then the following statements hold.
\begin{enumerate}

\item As $m\to\infty$,
$
\mathsf x_\ell^{(m)}\to(i,-i)
$
and
\begin{equation}\label{eq:L-limit}
        \frac{1}{m^2}\mathbf L_m\longrightarrow\mathbf L_\star,
\end{equation}
where
$
\mathbf L_\star\in\mathbb S^s
$
is a fixed real symmetric $s\times s$ matrix, independent of $m$.

\item There exist a constant $\gamma>0$ and an integer $m_0>0$ such that $
        \lambda_{\min}(\mathbf L_\star)<-\gamma<0. $
Let $c_\star\in\R^s$ be a corresponding real unit eigenvector.  Since
$\mathbf L_\star$ is independent of $m$, the same $c_\star$ is used for
every $m\ge m_0$.  

\item There exists an integer $m_1\ge m_0$ such that, for every
$m\ge m_1$, define
$
\widetilde c_m:=c_\star/m, 
$ then
\begin{equation}\label{eq:L-negative-limit}
 \widetilde c_m^*\mathbf L_m\widetilde c_m
 < -\frac{\gamma}{2}<0 .
\end{equation}

\item If $\widehat Q_{\rm av}\succeq0$, there exists an integer
$m_2\ge m_0$ such that, for every $m\ge m_2$,
\begin{equation}\label{eq:R-vanishes}
 \left|
 \widetilde c_m^*\mathbf R_m\widetilde c_m
 \right|
 <\frac{\gamma}{4},
\end{equation}
uniformly over all positive semidefinite normal-form Gram
representations.
\end{enumerate}
\end{proposition}

\begin{proof}
We prove the four statements in order. 
For the first statement, the entries of $\mathbf L_m$ are explicit finite
sums determined by $D,K,B$ and the points $\mathsf x_\ell^{(m)}$.  Since
$
\mathsf x_\ell^{(m)}\to(i,-i)
$
at rate $m^{-1/2}$, the corresponding expressions have a leading term
of order $m^2$.  For $1\le\ell,h\le s$, define
\[
 A_{\ell h}:=
 |a_\ell+\overline{a_h}|^2,
 \qquad
 C_{\ell h}:=
 4\,\Re(a_\ell)\Re(a_h).
\]
Expanding the leading term gives
\[
 \frac{1}{m^2}(\mathbf L_m)_{\ell h}
 \longrightarrow
 2\left(
 \frac{2}{A_{\ell h}^2}
 +\frac{1}{C_{\ell h}^2}
 -\frac{1}{A_{\ell h}C_{\ell h}}
 \right)
 =:(\mathbf L_\star)_{\ell h}.
\]
Thus
$
 \mathbf L_\star:=
 \left[
 2\left(
 \frac{2}{A_{\ell h}^2}
 +\frac{1}{C_{\ell h}^2}
 -\frac{1}{A_{\ell h}C_{\ell h}}
 \right)
 \right]_{\ell,h=1}^s
 \in\mathbb S^s
$
is a fixed real symmetric $s\times s$ matrix, independent of $m$.  Since
$s$ is fixed, entrywise convergence implies convergence in any matrix
norm, proving
$
        \frac{1}{m^2}\mathbf L_m\longrightarrow\mathbf L_\star.
$
% The complete calculation and finite-tail estimates are given in
% Appendix~\ref{app:quantitative}.

\noindent 
\textbf{For the second statement}, the fixed two-family choice of test points
determines $\mathbf L_\star$ completely.  Appendix~
\ref{app:negative-witness} gives an exact finite certificate showing that
$\mathbf L_\star$ has a strictly negative eigenvalue.  Hence there exists
$\gamma>0$ such that
\[
        \lambda_{\min}(\mathbf L_\star)<-\gamma<0.
\]
Let $c_\star\in\R^s$ be a corresponding real unit eigenvector:
$
 \mathbf L_\star c_\star
 =\lambda_{\min}(\mathbf L_\star)c_\star, 
 c_\star^*c_\star=1.
$
The matrix $\mathbf L_\star$, the constant $\gamma$, and the vector
$c_\star$ are fixed and independent of $m$.  Choose $m_0$ sufficiently
large that the asymptotic estimates below are valid, and define
$
\widetilde c_m:=c_\star/m
$
for every $m\ge m_0$.

\noindent
\textbf{For the third statement,}
\[
 \widetilde c_m^*\mathbf L_m\widetilde c_m
 =
 c_\star^*
 \left(\frac{\mathbf L_m}{m^2}\right)c_\star
 \longrightarrow
 c_\star^*\mathbf L_\star c_\star
 =
 \lambda_{\min}(\mathbf L_\star)
 <-\gamma.
\]
Therefore there exists an integer $m_1\ge m_0$ such that, for every
$m\ge m_1$,
$
 \widetilde c_m^*\mathbf L_m\widetilde c_m
 < -\frac{\gamma}{2}<0.
$

\noindent
\textbf{For the fourth statement}, suppose that
$\widehat Q_{\rm av}\succeq0$.  Positivity of the mixed block first
forces the low-order mixed directions into its kernel.  The resulting
improvement at the conjugate test points can then be continued to the
realigned points appearing in $\mathbf R_m$.  This yields
\[
 \sup_{\widehat Q_{\rm av}\succeq0}
 \left|
 c_\star^*\frac{\mathbf R_m}{m^2}c_\star
 \right|
 \longrightarrow0.
\]
The detailed argument is given in Appendix~\ref{app:quantitative}.
Hence there exists an integer $m_2\ge m_0$ such that, for every
$m\ge m_2$ and every positive semidefinite normal-form Gram
representation,
$
 \left|
 \widetilde c_m^*\mathbf R_m\widetilde c_m
 \right|
 =
 \left|
 c_\star^*\frac{\mathbf R_m}{m^2}c_\star
 \right|
 <\frac{\gamma}{4}.
$
This completes the proof.
\end{proof}

% \begin{proof}
% We prove the four statements in order. 
% For the first statement, the entries of $\mathbf L_m$ are explicit finite
% sums determined by $D,K,B$ and the fixed test points
% $\mathsf x_\ell^{(m)}$.  Since
% $
% \mathsf x_\ell^{(m)}\to(i,-i)
% $
% at rate $m^{-1/2}$, put, for $1\le\ell,h\le s$,
% \[
%  A_{\ell h}:=
%  |a_\ell+\overline{a_h}|^2,
%  \qquad
%  C_{\ell h}:=
%  4\,\Re(a_\ell)\Re(a_h).
% \]
% Using
% $
% 1-e^{-z/\sqrt m}=z/\sqrt m+O(m^{-1}),
% $
% the leading terms of the sums defining $\mathbf L_m$ give
% \[
%  \frac{1}{m^2}(\mathbf L_m)_{\ell h}
%  \longrightarrow
%  2\left(
%  \frac{2}{A_{\ell h}^2}
%  +\frac{1}{C_{\ell h}^2}
%  -\frac{1}{A_{\ell h}C_{\ell h}}
%  \right)
%  =:(\mathbf L_\star)_{\ell h}.
% \]
% Since $s$ is fixed,
% entrywise convergence implies convergence in any matrix norm.  Hence
% \[
%         \frac{1}{m^2}\mathbf L_m\longrightarrow\mathbf L_\star,
% \]
% where $\mathbf L_\star\in\mathbb S^s$ is a fixed real symmetric matrix,
% independent of $m$.
% % \footnote{More details of the the tail estimate is given in \Cref{B:lem:tail}}.

% \noindent \textbf{For the second statement,} $\mathbf L_\star$ is completely determined by
% the fixed two-family set of test speeds.  Appendix~
% \ref{app:negative-witness} gives an exact rational certificate showing
% that
% \[
%         \lambda_{\min}(\mathbf L_\star)<-2^{-18}<0.
% \]
% Thus, for example, we may take
% $
% \gamma=2^{-18}.
% $
% Let $c_\star\in\R^s$ be a corresponding real unit eigenvector,
% \[
%  \mathbf L_\star c_\star
%  =\lambda_{\min}(\mathbf L_\star)c_\star,
%  \qquad
%  c_\star^*c_\star=1.
% \]
% The matrix $\mathbf L_\star$, the constant $\gamma$, and the vector
% $c_\star$ are fixed and independent of $m$.  Choose an integer $m_0$
% large enough that the asymptotic estimates below apply, and define
% $
% \widetilde c_m:=c_\star/m
% $
% for $m\ge m_0$.

% \noindent \textbf{For the third statement,}
% \[
%  \widetilde c_m^*\mathbf L_m\widetilde c_m
%  =
%  c_\star^*
%  \left(\frac{\mathbf L_m}{m^2}\right)c_\star
%  \longrightarrow
%  c_\star^*\mathbf L_\star c_\star
%  =
%  \lambda_{\min}(\mathbf L_\star)
%  <-\gamma.
% \]
% Hence there exists an integer $m_1\ge m_0$ such that, for every
% $m\ge m_1$,
% \[
%  \widetilde c_m^*\mathbf L_m\widetilde c_m
%  < -\frac{\gamma}{2}<0.
% \]

% \noindent \textbf{For the fourth statement}, suppose that
% $\widehat Q_{\rm av}\succeq0$.  The low-total mixed directions are then
% forced into the kernel of every feasible mixed block.  This gives an
% improved estimate at the conjugate test points.  The qualitative argument
% in \Cref{B:prop:R} gives
% \[
%  \sup_{\widehat Q_{\rm av}\succeq0}
%  \left|
%  c_\star^*\frac{\mathbf R_m}{m^2}c_\star
%  \right|
%  \longrightarrow0.
% \]
% Therefore there exists an integer $m_2\ge m_0$ such that, for every
% $m\ge m_2$ and every positive semidefinite normal-form Gram
% representation,
% $
%  \left|
%  \widetilde c_m^*\mathbf R_m\widetilde c_m
%  \right|
%  =
%  \left|
%  c_\star^*\frac{\mathbf R_m}{m^2}c_\star
%  \right|
%  <\frac{\gamma}{4}.
% $
% This completes the proof.
% \end{proof}

Finally, we piece together the preceding lemmas to conclude the section and prove Theorem~\ref{thm:negative}.

\begin{proof}[Proof of Theorem~\ref{thm:negative}]
Let
$
M_0:=\max\{m_1,m_2\}
$
and suppose, for contradiction, that $\mathcal F_m$ is SoS for some
$m\ge M_0$.  By Lemma~\ref{lem:finite-test}, the corresponding finite
Gram test satisfies
$
\mathcal G_m\succeq0,
$
and therefore
$
        \widetilde c_m^*\mathcal G_m\widetilde c_m\ge0.
$
On the other hand, \Cref{lem:finite-decomposition,lem:negative-limit} give
\[
\begin{aligned}
\widetilde c_m^*\mathcal G_m\widetilde c_m
&\le
\widetilde c_m^*\mathbf L_m\widetilde c_m
+
\left|
\widetilde c_m^*\mathbf R_m\widetilde c_m
\right| <
-\frac{\gamma}{2}+\frac{\gamma}{4}
=
-\frac{\gamma}{4}<0,
\end{aligned}
\]
contradicting $\mathcal G_m\succeq0$. Therefore $\mathcal F_m$ is not SoS for every sufficiently large corner
depth $m$.  In particular, choose any integer $m_\star\ge M_0$.
By Lemma~\ref{lem:corner}, $\mathcal F_{m_\star}$ is a restriction of
$F_N$ for every $N\ge2m_\star$.  Hence $F_N$ is not SoS for every
$N\ge2m_\star$.  Taking
$
N_0:=2m_\star
$
proves the theorem.
\end{proof}

\section{Subclasses of Toeplitz BW Form that are SoS}
\label{sec:positive-subclasses}

The failure of SoS representability at large orders does not persist
under every structural restriction.  In this section, we show that 
the Toeplitz B\"ottcher--Wenzel form is SoS at every order whenever
one factor is symmetric or skew-symmetric. 
The other factor remains an arbitrary real Toeplitz matrix.


The proof proceeds in two steps.  First, \Cref{lem:positive-split} shows
that the Toeplitz B\"ottcher--Wenzel form with one symmetric factor can
be split into a term involving two symmetric Toeplitz matrices and a
term involving a symmetric and a skew-symmetric matrix.
\Cref{prop:same-symmetry} proves that the first term is SoS, while
\Cref{lem:positive-mixed} proves the same for the second.  The
skew-symmetric case follows analogously, with the roles of the symmetric
and skew-symmetric parts interchanged. 
After completing the proof, we discuss the exterior-product
interpretation, and its relation to the two-matrix De Smet--Dillen--Verstraelen--Vrancken (DDVV) inequality. 
For notational simplicity, throughout this section we write $F(X,Y)$
for $F_N(x,y)$ in \eqref{eq:FN}, where $X$ and $Y$ are the 
matrices of size $N$.  

% For real matrices $U,V$ of the same order, write
% \[
%  F(U,V)
%  :=
%  2\norm{U}_F^2\norm{V}_F^2
%  -2\langle U,V\rangle_F^2
%  -\norm{[U,V]}_F^2.
% \]
% Thus $F_N(x,y)=F(X,Y)$ and $F(U,V)=F(V,U)$.


\paragraph{Splitting the factor}
\label{subsec:positive-split}
Clearly, every real Toeplitz matrix $X$ admits an orthogonal
decomposition into a symmetric Toeplitz matrix and a skew-symmetric
Toeplitz matrix:
\begin{equation}\label{eq:positive-split}
 X=S_x+K_x,
 \qquad
 S_x=\frac{X+X^\top}{2},
 \qquad
 K_x=\frac{X-X^\top}{2}.
\end{equation}
Here $S_x$ is symmetric Toeplitz and $K_x$ is skew-symmetric Toeplitz.
Moreover, they are orthogonal with respect to the Frobenius inner
product,
$
 \langle S_x,K_x\rangle_F=0,
$
and hence
$
\|X\|_F^2=\|S_x\|_F^2+\|K_x\|_F^2.
$
The following lemma shows that this decomposition also splits the
B\"ottcher--Wenzel form.

\begin{lemma}[The Splitting Lemma]
\label{lem:positive-split}
Let $X$ be a real Toeplitz matrix, and let $S$ and $K$ be real
symmetric and skew-symmetric matrices, respectively, of the same order.
Then
\begin{equation}\label{eq:positive-split-identities}
 \begin{aligned}
 F(X,S)=F(S,X)&=F(S,S_x)+F(S,K_x),\\
 F(X,K)=F(K,X)&=F(K,S_x)+F(K,K_x),
 \end{aligned}
\end{equation}
where $S_x$ and $K_x$ are defined in \eqref{eq:positive-split}.
\end{lemma}

\begin{proof}
Symmetric and skew-symmetric matrices are orthogonal with respect to
the Frobenius inner product.  In particular,
\[
 \langle S,K_x\rangle_F=0,
 \qquad
 \|S_x+K_x\|_F^2
 =
 \|S_x\|_F^2+\|K_x\|_F^2.
\]
Moreover, $[S,S_x]$ is skew-symmetric, whereas $[S,K_x]$ is symmetric.
Hence
\[
 \|[S,S_x+K_x]\|_F^2
 =
 \|[S,S_x]\|_F^2+\|[S,K_x]\|_F^2.
\]
Substitution into the definition of $F$ proves the first identity.
For the second identity, use
$
\langle K,S_x\rangle_F=0
$
and the fact that $[K,S_x]$ is symmetric, while $[K,K_x]$ is
skew-symmetric.  The same argument then gives the second identity.
\end{proof}

\paragraph{The two types of pairs}
\label{subsec:positive-pairs}

The splitting lemma \Cref{lem:positive-split} reduces the theorem to two
types of BW forms.  \Cref{lem:positive-mixed} shows that the
symmetric--skew-symmetric term in \eqref{eq:positive-split-identities}
is SoS.  In fact, this result holds for any symmetric--skew-symmetric
pair, even without the Toeplitz assumption.  \Cref{prop:same-symmetry}
shows that the same-symmetry term in
\eqref{eq:positive-split-identities} is also SoS.

\begin{lemma}[Symmetric--skew-symmetric pair]
\label{lem:positive-mixed}
For every real symmetric matrix
$S=(s_{ij})_{1\le i,j\le N}$ and real skew-symmetric matrix
$K=(k_{ij})_{1\le i,j\le N}$ of order $N$,
\begin{equation}\label{eq:positive-mixed}
 F(S,K)
 =
 \norm{SK+KS}_F^2
 +
 4\sum_{\ell=1}^N
 \sum_{1\le i<j<k\le N}
 (s_{i\ell}k_{jk}-s_{j\ell}k_{ik}+s_{k\ell}k_{ij})^2.
\end{equation}
\end{lemma}

\begin{proof}

Since $\langle S,K\rangle_F=0$,
$
 F(S,K)=2\norm{S}_F^2\norm{K}_F^2-\norm{SK-KS}_F^2.
$
For a vector $v\in\mathbb R^N$, the exterior-product Lagrange
identity gives
\[
 \frac12\norm{v}^2\norm{K}_F^2-\norm{Kv}^2
 =
 \sum_{i<j<k}
 (v_i k_{jk}-v_j k_{ik}+v_k k_{ij})^2.
\]
This identity follows directly by expansion.
Apply it to each column of $S$ and sum to obtain
\[
 2\norm{S}_F^2\norm{K}_F^2
 =
 4\norm{KS}_F^2
 +
 4\sum_{\ell=1}^N\sum_{i<j<k}
 (s_{i\ell}k_{jk}-s_{j\ell}k_{ik}+s_{k\ell}k_{ij})^2.
\]
Finally, $(SK)^\top=-KS$, so
$\norm{SK}_F=\norm{KS}_F$.
The parallelogram identity therefore yields
$
 4\norm{KS}_F^2-\norm{SK-KS}_F^2
 =
 \norm{SK+KS}_F^2.
$
Combining the two identities proves
\eqref{eq:positive-mixed}.
Every entry of $SK+KS$, and every expression inside the
remaining squares, is a homogeneous quadratic polynomial
in the independent entries of $S$ and $K$.
\end{proof}

\begin{proposition}[Both symmetric or skew-symmetric pair]
\label{prop:same-symmetry}
For every $N\ge2$, $F(X,Y)$ is SoS if $X,Y$ are both symmetric
Toeplitz or both skew-symmetric Toeplitz.
\end{proposition}

\begin{proof}[Proof sketch]

In either case, the matrices are determined by their nonnegative
diagonal parameters. We use the wedge coordinates introduced in
\Cref{subsec:wedge-coordinates}. Set $m=N-1$ and reverse the nonzero
diagonal labels by writing
\[
 \hat z_{ab}=z_{N-b,N-a},
 \qquad 1\le a<b\le m.
\]
This indexing replaces the diagonal-length weights
$(N-a)(N-b)$ by $ab$.
We rewrite the Toeplitz form as
\begin{equation}\label{eq:positive-reduction}
 \begin{aligned}
 F^{\rm sym}(X,Y)
 =
 4N\sum_{b=1}^{N-1}(N-b)z_{0b}^2
 +4\Psi_{N-1}(\hat z),\qquad
 F^{\rm skew}(X,Y)
 =4\Psi_{N-1}(\hat z).
 \end{aligned}
\end{equation}
Note that the full calculation is given in
\Cref{app:same-symmetry-reduction}.
Here $\Psi_m(\hat z)$ is defined by
\begin{equation}\label{eq:positive-reduced-form}
 \Psi_m(\hat z)
 :=
 2\sum_{1\le a<b\le m}ab\,\hat z_{ab}^2
 -
 \sum_{\delta=1}^{m-1}
 \sum_{\substack{k\ge0\\2k<m-\delta}}
 \left(
   \sum_{p=k+1}^{m-\delta-k}\hat z_{p,p+\delta}
 \right)^2.
\end{equation}
Intuitively, the inner sum in
\eqref{eq:positive-reduced-form} selects a window from the sequence
$
\hat z_{1,1+\delta},\ldots,\hat z_{m-\delta,m},
$
successively removing $k$ terms from each end. 
All we need to prove next is that
$\Psi_m(\hat z(x,y))$ is SoS. For example, when $m=4$ and $\delta=1$,
the two windows give
\[
 (\hat z_{12}+\hat z_{23}+\hat z_{34})^2
 \qquad\text{and}\qquad
 \hat z_{23}^2.
\]
The essential observation is that the wedges satisfy the Pl\"ucker
relations recalled in \Cref{subsec:wedge-coordinates}. These relations
allow the cross terms in the negative window squares to be re-paired
without changing the polynomial in $x,y$. For instance,
\[
 -2\hat z_{12}\hat z_{34}
 =
 -2\hat z_{13}\hat z_{24}
 +2\hat z_{14}\hat z_{23}.
\]
Applying the prescribed re-pairing rule gives three types of
squares: window sums along a fixed index total, differences of
wedges with the same gap, and individual wedges.
\Cref{prop:positive-window-certificate} gives the
complete identity and its coefficients. The window and
difference terms have nonnegative integer weights.
After allocating the diagonal terms to these squares, the
remaining coefficient of $\hat z_{p,p+\delta}^2$, denoted
$s_\delta(p)$ in \Cref{app:positive-windows}, satisfies
$
 s_\delta(p)\ge p^2>0.
$
This estimate ensures that completing the squares leaves no
negative remainder.

Consequently, $\Psi_m$ is a sum of squares of linear combinations
of the wedges, after using their Pl\"ucker relations.
Since each wedge is quadratic in $x,y$, this is an SoS
representation in the original variables.
Equation~\eqref{eq:positive-reduction} proves both assertions.
\end{proof}


Combining the splitting in \Cref{lem:positive-split} with the two
results above gives the following all-order positive result.

\begin{theorem}[One symmetric or skew-symmetric factor]
\label{thm:one-sided}
For every $N\ge2$, the Toeplitz B\"ottcher--Wenzel form $F(X,Y)$ is SoS
whenever either $X$ or $Y$ is symmetric or skew-symmetric Toeplitz.
The other factor may be an arbitrary real Toeplitz matrix.
\end{theorem}

\begin{proof}
Suppose first that $X$ is symmetric Toeplitz.  By
\Cref{lem:positive-split},
$
 F(X,Y)=F(X,S_y)+F(X,K_y),
$
where the two terms are SoS by \Cref{prop:same-symmetry} and
\Cref{lem:positive-mixed}, respectively.  The skew-symmetric case is
identical with the two results interchanged.  Finally,
$F(X,Y)=F(Y,X)$ gives the same conclusion when the structural
assumption is imposed on $Y$.
\end{proof}

\begin{remark}
The splitting above relies on one factor having a fixed symmetry.  If
both factors are decomposed into symmetric and skew-symmetric parts,
additional cross terms appear and the four resulting BW forms no longer
give the whole expression.  Thus the argument proves the one-sided
result above, but does not extend directly to unrestricted Toeplitz
pairs.  The full expansion is recorded in
Appendix~\ref{app:positive-coupling}.
\end{remark}

\paragraph{An exterior-product regularization viewpoint}
The positive part of the B\"ottcher--Wenzel form admits a useful
regularization interpretation.  By the weighted Lagrange identity,
\[
 \|X\|_F^2\|Y\|_F^2-\langle X,Y\rangle_F^2
 =
 z^\top W_Nz,
 \qquad W_N\succ0,
\]
where $z=x\wedge y$ denotes the wedge coordinates introduced in
\Cref{subsec:wedge-coordinates}.  Since the commutator is linear in
$z$, say $[X,Y]=\mathscr C_Nz$, the generalized form
\[
 F_\tau(X,Y)
 :=
 \tau\bigl(
 \|X\|_F^2\|Y\|_F^2-\langle X,Y\rangle_F^2
 \bigr)
 -
 \|[X,Y]\|_F^2
\]
becomes
\[
 F_\tau(X,Y)
 =
 z^\top
 \bigl(\tau W_N-\mathscr C_N^*\mathscr C_N\bigr)z.
\]
Thus, for each fixed $N$, the exterior-product term acts as a
positive-definite regularizer in wedge space, and sufficiently large
$\tau$ makes the lifted quadratic form positive semidefinite.
\rev{Related interactions between regularization, global optimality, and
SoS exactness for quartic polynomial models are studied in
\cite{zhuCartis2026regularizedSOS,zhuCartis2025global,zhuCartis2024sosTaylor}.}
The B\"ottcher--Wenzel form corresponds to the fixed value $\tau=2$. 
This differs from a radial regularization such as
$\sigma\|(x,y)\|^4$: the exterior-product term vanishes whenever
$X$ and $Y$ are linearly dependent, and is therefore degenerate in
the original variables even though it is positive definite in wedge
coordinates.
\paragraph{Relation to the two-matrix DDVV form}

A closely related commutator polynomial arises from the
De Smet--Dillen--Verstraelen--Vrancken (DDVV) inequality for real
symmetric matrices \cite{lu2017,ge2024published}.  For two such
matrices, define
\[
 \mathcal D_N(X,Y)
 :=
 \bigl(\norm{X}_F^2+\norm{Y}_F^2\bigr)^2
 -2\norm{[X,Y]}_F^2.
\]
A direct expansion gives
\begin{equation}\label{eq:positive-ddvv-relation}
 \mathcal D_N(X,Y)
 =
 2F(X,Y)
 +\bigl(\norm{X}_F^2-\norm{Y}_F^2\bigr)^2
 +4\langle X,Y\rangle_F^2.
\end{equation}
Hence an SoS representation of $F$ immediately gives one for
$\mathcal D_N$.  The converse does not follow, since the last two
terms in \eqref{eq:positive-ddvv-relation} are additional squares. 
The two-matrix DDVV form itself admits an SoS representation for
arbitrary real symmetric matrices.  Thus the DDVV result does not
replace the Toeplitz argument in \Cref{prop:same-symmetry}; rather,
\eqref{eq:positive-ddvv-relation} shows precisely how the two
quartic forms differ.
\section{Exact certificates for the Toeplitz BW form at small orders}
\label{sec:exact-certificates}

We have shown that, for sufficiently large $N$, the general Toeplitz
B\"ottcher--Wenzel form is not SoS.  In the previous section, we also
showed that several structured subclasses remain SoS at every order.

This naturally raises the question of what happens for the small orders
that arise in applications, when no additional symmetry structure is
assumed.  Is the unrestricted Toeplitz BW form still SoS?
For small orders, the answer is positive.  In this section, we show
that the unrestricted Toeplitz BW form is SoS for every
$
2\le N\le50.
$
Moreover, for
$
2\le N\le20,
$
we provide a Lean-verified proof of the SoS property\footnote{Note that the orders $21$--$50$ are certified by exact arithmetic but are
not included in the Lean verification.}.

Together with the large-order negative result, this leaves a substantial
gap.  The quantitative argument in Appendix~\ref{app:quantitative}
gives the explicit sufficient bound
$
N_0=\SufficientOrder
$
for non-SoS representability, whereas the exact positive construction
below has currently been verified only through order $50$.  The latter
stops because the positive increment used in the induction has only
been constructed and verified through the step $n=49$; this does not
imply that $F_{51}$ is not SoS.  Thus the exact SoS status of the
unrestricted Toeplitz BW form for
$
50<N<N_0
$
remains open.  Determining the first non-SoS order, or substantially
narrowing this gap, is a natural direction for future work.

\paragraph{Toeplitz BW SoS for $2\le N\le50$}\label{subsec:certificate-chain}
We establish the result by an exact induction. Starting from $F_2$,
each step expresses the next Toeplitz BW form as a positive multiple
of a linear substitution of the preceding form, plus an SoS increment.

Here $N$ denotes the matrix order in the final result, while $n$ is
used only as the induction index. At the step from $n$ to $n+1$, let
$z^{(n+1)}$ denote the wedge vector associated with the symbol variables
of $F_{n+1}$.

\begin{lemma}[Exact positive increments]
\label{lem:finite-increments}
For each $n=2,\ldots,49$, the supplied rational data specify a
positive rational number $\alpha_n$, a rational linear map
$\psi_n:\R^{4n+2}\to\R^{4n-2}$, and a rational symmetric matrix
$\mathbf H_n\succeq0$ such that
\begin{equation}\label{eq:finite-chain}
 F_{n+1}
 =
 \alpha_n(F_n\circ\psi_n)
 +(z^{(n+1)})^\top\mathbf H_nz^{(n+1)}.
\end{equation}
\end{lemma}

\begin{proof}
Each identity is verified by reconstructing the two quartics from
their defining Toeplitz matrices and comparing all polynomial
coefficients over $\mathbb Q$. Positive semidefiniteness of
$\mathbf H_n$ is checked by exact rational matrix calculations after
removing explicitly verified kernel directions. The data for all
$48$ steps and the verification procedure are given in
Appendix~\ref{app:finite-certificates}.
\end{proof}


\begin{theorem}[The finite positive range]
\label{thm:finite-positive}
For every integer $2\le N\le50$, $F_N$ is a sum of squares of real
homogeneous quadratic polynomials. For $2\le N\le20$, the same
statement is also formally verified in Lean~4.
\end{theorem}

\begin{proof}
At order two, direct expansion gives
$
 F_2=4z_{-1,0}^2+4z_{0,1}^2.
$
Suppose $F_n=\sum_j q_j^2$, where the $q_j$ are real homogeneous
quadratics. Since $\mathbf H_n\succeq0$, it admits a real factorization
$\mathbf H_n=V_n^\top V_n$. Equation~\eqref{eq:finite-chain} gives
\[
 F_{n+1}
 =
 \sum_j\bigl(\sqrt{\alpha_n}\,q_j\circ\psi_n\bigr)^2
 +
 \sum_i\bigl(V_nz^{(n+1)}\bigr)_i^2.
\]
Each factor is again a homogeneous quadratic. Induction through
$n=49$ therefore proves the result for every $2\le N\le50$.
\end{proof}

The substitutions in \eqref{eq:finite-chain} allow the SoS form at one
order to be reused at the next. For $2\le n\le16$, a simple rescaling
of $F_n$ is sufficient. For $17\le n\le49$, the substitution also
damps the symbol at offset $a$ by a rational factor $q_n^{|a|}$,
with $0<q_n\le1$, so that the remaining increment admits a positive
semidefinite Gram representation. The precise constructions are given
in Appendix~\ref{app:finite-certificates}. 

\paragraph{Lean verification through order $20$.} For $2\le N\le20$, the SoS assertion is also formally verified in
Lean~4.  In particular, the range $13\le N\le20$ extends the previously
known positive range through order $12$.  The Lean formalization works
with the original Toeplitz BW polynomial and proves directly that it is
a finite sum of real homogeneous quadratic squares.  The corresponding
formal statement and source dependencies are recorded in
Appendix~\ref{app:positive-lean}.

This verification is distinct from the exact rational induction in
\Cref{subsec:certificate-chain}.  The rational argument proves the SoS
property for every $2\le N\le50$ by checking the identities
\eqref{eq:finite-chain} and the positive semidefiniteness of the
associated rational matrices.  The Lean proof, on the other hand,
provides an independent formal verification for each order
$2\le N\le20$ and need not use the same Gram representations. 
Thus,
\[
 \begin{array}{c|c}
 \text{orders}&\text{verification}\\\hline
 2\le N\le20&\text{exact rational verification and Lean~4}\\
 21\le N\le50&\text{exact rational verification}
 \end{array}
\]
Both results concern unrestricted real Toeplitz pairs.  The status of
the unrestricted form remains open for the intermediate range
$
50<N<N_0.
$
\section{Conclusion}\label{sec:conclusion}

This paper shows that Toeplitz structure alone is not sufficient to
guarantee a sum-of-squares representation of the
B\"ottcher--Wenzel form.  The main negative result is obtained by
restricting the Toeplitz form to a corner-supported family whose
structure is independent of the matrix order.  After describing the
complete family of Gram representations of this restricted form, we
separate a fixed Gram-invariant contribution from the remaining
Gram-dependent part.  A finite complex test then detects a negative
limiting direction in the former, while the latter vanishes uniformly
over all positive semidefinite Gram representations.  This gives the
contradiction and proves that the unrestricted Toeplitz BW form is not
SoS for every sufficiently large order.

At the same time, the failure is not universal across structured
Toeplitz classes.  We prove that if either factor is symmetric or
skew-symmetric Toeplitz, then the BW form remains SoS at every order,
with the other factor completely unrestricted within the real Toeplitz
class.  The proof follows from an orthogonal decomposition into
symmetric and skew-symmetric parts, together with SoS identities for
pairs of the same symmetry and for symmetric--skew-symmetric pairs.
Thus the negative result is caused not by Toeplitz structure itself,
but by the interaction between the two symmetry components when both
factors contain them.

These two results leave several natural questions open.  Most
immediately, the smallest order at which the unrestricted Toeplitz BW
form ceases to be SoS is still unknown.  The present negative argument
gives only a sufficient large-order threshold, while exact computations
show that the form remains SoS throughout a substantial initial range.
A second direction is to identify further structured subspaces, beyond
the symmetric and skew-symmetric cases considered here, on which the
interaction between the two symmetry components can still be controlled.
More generally, it would be interesting to understand whether the
corner and Gram-realignment mechanisms developed here can be used to
locate sharper boundaries between nonnegativity and SoS in other
structured polynomial matrix inequalities.

\paragraph{Declaration of generative AI and AI-assisted technologies}
During the preparation of this work the authors used large language model
systems (Claude, Anthropic; Codex models, OpenAI) to generate candidate
constructions and code, to run computations, to draft mathematical
arguments, and to improve the readability of the text. After using these
tools, the authors reviewed and edited the content as needed and take full
responsibility for the content of the publication.

\paragraph{Data availability}
\rev{The exact certificate data, verification scripts, and Lean~4
formalization are available at
\url{https://github.com/erenup/toeplitz-bw-not-sos}.  The archived release
used for the results reported here is tag \texttt{v1.0}; the release
documentation is recorded at commit \texttt{7126c08}.}


\appendix

\section{Proof of the general Gram normal form}
\label{app:normal-form-proof}

This appendix proves Lemma~\ref{lem:normal-form}.  We follow the three
statements of the lemma in order.

\begin{proof}[Proof of Lemma~\ref{lem:normal-form}]
Put $r=\binom m2$.

\medskip
\noindent
\textbf{1. Reduction of the Gram family.}
By Lemma~\ref{lem:frame}, every Gram representation of $\calF_m$
differs from the fixed representation
\[
G_m=
\begin{pmatrix}
D&K&0\\
K&D&0\\
0&0&B
\end{pmatrix}
\]
by Pl\"ucker corrections.  We classify these corrections according to
how their four labels are distributed between the two sides of the
Toeplitz corner.

Four labels on the same side produce corrections acting entirely within
one pure wedge family.  Such a correction has entries
$
(\mathcal E_0)_{ab,cd}=\Omega_{abcd},
$
where $\Omega$ is totally antisymmetric in its four indices.  These are
the same-side Pl\"ucker corrections, or equivalently the pure
four-forms.

Two labels on each side produce the Pl\"ucker relations that couple the
two pure families with the mixed family.  After the pure cross
correction is collected into a symmetric matrix
$
\mathcal E_1\in\mathbb S^r,
$
the corresponding change in the mixed block is determined by the same
Pl\"ucker relations.  We denote this induced matrix by
$
\calR\mathcal E_1\in\mathbb S^{m^2}.
$
Equivalently, after antisymmetrically extending the entries of
$\mathcal E_1$ in each pure index pair, the realignment map is given by
\[
 (\calR\mathcal E_1)_{ac,bd}
 =(\mathcal E_1)_{ab,cd},
 \qquad
 (\calR\mathcal E_1)_{ad,bc}
 =-(\mathcal E_1)_{ab,cd}.
\]
Thus $\calR$ simply re-pairs the four indices from the two pure wedge
pairs into two mixed wedge pairs.

Corrections involving three labels on one side and one label on the
other produce pure--mixed blocks.  To remove these terms, change the
sign of all variables on the negative side.  This transformation leaves
the pure coordinates $u,v$ unchanged and changes the sign of every
mixed coordinate $w$.  The polynomial $\calF_m$ is unchanged, so if
$Q\succeq0$ is a Gram representation, averaging $Q$ with its image
under this symmetry preserves both the represented polynomial and
positive semidefiniteness, while eliminating all pure--mixed blocks.

Next interchange the two sides of the Toeplitz corner.  Under this
symmetry,
\[
                         (u,v)\longmapsto(-v,-u),
 \qquad
                         w_{ab}\longmapsto-w_{ba}.
\]
A second averaging therefore makes the two pure diagonal corrections
equal and makes the pure cross correction symmetric.  Consequently,
whenever a positive semidefinite Gram representation exists, there is
one of the form
\[
 Q_{\rm av}
 =
 G_m+
 \begin{pmatrix}
 \mathcal E_0&\mathcal E_1&0\\
 \mathcal E_1&\mathcal E_0&0\\
 0&0&\calR\mathcal E_1
 \end{pmatrix},
\]
which is exactly \eqref{eq:general-gram-1}.  Both averaging operations
are convex averages of positive semidefinite matrices and therefore
preserve positive semidefiniteness.

\medskip
\noindent
\textbf{2. Block diagonalization.}
Write
\[
 A:=D+\mathcal E_0,
 \qquad
 C:=K+\mathcal E_1.
\]
Then the two pure blocks of $Q_{\rm av}$ have the form
\[
 \begin{pmatrix}
 A&C\\
 C&A
 \end{pmatrix}.
\]
Consider the orthogonal matrix
\[
 U=
 \begin{pmatrix}
 \frac1{\sqrt2}I_r&\frac1{\sqrt2}I_r&0\\
 \frac1{\sqrt2}I_r&-\frac1{\sqrt2}I_r&0\\
 0&0&I_{m^2}
 \end{pmatrix}.
\]
It corresponds to the change of coordinates
\[
 (u,v,w)
 \longmapsto
 \left(
 \frac{u+v}{\sqrt2},
 \frac{u-v}{\sqrt2},
 w
 \right).
\]
Since $U^\top U=I$, direct multiplication gives
\[
\begin{aligned}
 U^\top Q_{\rm av}U
 &=
 \begin{pmatrix}
 A+C&0&0\\
 0&A-C&0\\
 0&0&B+\calR\mathcal E_1
 \end{pmatrix}                                                   \\[1mm]
 &=
 \begin{pmatrix}
 D+\mathcal E_0+K+\mathcal E_1&0&0\\
 0&D+\mathcal E_0-K-\mathcal E_1&0\\
 0&0&B+\calR\mathcal E_1
 \end{pmatrix}.
\end{aligned}
\]
Hence
\[
 P_+=D+\mathcal E_0+K+\mathcal E_1,\qquad
 P_-=D+\mathcal E_0-K-\mathcal E_1,\qquad
 M=B+\calR\mathcal E_1,
\]
which proves \eqref{eq:general-gram-2}--\eqref{eq:sectors}.

\medskip
\noindent
\textbf{3. Equivalence with the SoS condition.}
Suppose first that $\calF_m$ is SoS.  By Lemma~\ref{lem:frame}, it has a
positive semidefinite Gram representation in the restricted wedge
coordinates.  Part~1 shows that this representation can be replaced,
without losing positive semidefiniteness, by one of the form
\eqref{eq:general-gram-1}.  Thus there exist
$\mathcal E_0,\mathcal E_1$ for which $Q_{\rm av}\succeq0$.

Conversely, if such matrices $\mathcal E_0,\mathcal E_1$ exist and
$Q_{\rm av}\succeq0$, then
$
Q_{\rm av}=R^\top R
$
for some real matrix $R$.  Since $Q_{\rm av}$ is a Gram representation
of $\calF_m$,
\[
 \calF_m
 =z_r^\top Q_{\rm av}z_r
 =\|Rz_r\|^2,
\]
which is a sum of squares of homogeneous quadratic polynomials.
Therefore
\[
 \calF_m\text{ is SoS}
 \quad\Longleftrightarrow\quad
 \exists\,\mathcal E_0,\mathcal E_1
 \text{ such that }Q_{\rm av}\succeq0.
\]

Finally, $U$ is orthogonal, so
\[
 Q_{\rm av}\succeq0
 \quad\Longleftrightarrow\quad
 U^\top Q_{\rm av}U\succeq0.
\]
By the block diagonal form in Part~2, this is equivalent to
\[
                         P_+\succeq0,\qquad
                         P_-\succeq0,\qquad
                         M\succeq0.
\]
This proves the lemma.
\end{proof}
\section{Qualitative estimates for the finite Gram test}
\label{app:quantitative}

This appendix supplies the qualitative estimates used in
\Cref{lem:negative-limit}.  The main theorem only requires their
large-$m$ limits.  We also record, at the end of the appendix, the
quantitative constants used to obtain the explicit sufficient order in
\Cref{app:explicit-bound}.

\subsection{The scaled limit of the Gram-invariant part}

Let
\[
 \delta_m:=m^{-1/2},\qquad
 \mathsf x_\ell^{(m)}
 =(\xi_\ell^{(m)},\eta_\ell^{(m)})
 =\left(ie^{-\delta_m a_\ell},
        -ie^{-\delta_m\overline{a_\ell}}\right),
 \qquad 1\le\ell\le s.
\]
The speed parameters $a_1,\ldots,a_s$ are fixed throughout, with
$\Re a_\ell>0$.  Let $\mathbf L_{\infty,m}$ denote the matrix obtained
from $\mathbf L_m$ by extending the depth indices in $D,K,B$ from
$\{0,\ldots,m-1\}$ to all nonnegative integers.

For the pair $(\ell,h)$, set
\[
\begin{aligned}
 A_m^{\ell h}
 &:=(1-\xi_\ell^{(m)}\overline{\xi_h^{(m)}})
    (1-\eta_\ell^{(m)}\overline{\eta_h^{(m)}}),\\
 C_m^{\ell h}
 &:=(1-\xi_\ell^{(m)}\eta_\ell^{(m)})
    (1-\overline{\xi_h^{(m)}}\,\overline{\eta_h^{(m)}}),\\
 \Delta_m^{\ell h}
 &:=(1-\xi_\ell^{(m)}\overline{\eta_h^{(m)}})
    (1-\eta_\ell^{(m)}\overline{\xi_h^{(m)}}).
\end{aligned}
\]
Direct summation of the geometric series associated with $D,K,B$ gives
\begin{equation}\label{B:eq:Linfty}
 (\mathbf L_{\infty,m})_{\ell h}
 =
 2\left(
 \frac{2}{(A_m^{\ell h})^2}
 +\frac{1}{(C_m^{\ell h})^2}
 -\frac{1}{A_m^{\ell h}C_m^{\ell h}}
 -\frac{2}{(\Delta_m^{\ell h})^2}
 \right).
\end{equation}

\begin{lemma}[Finite-depth tail]\label{B:lem:tail}
For the fixed finite set of test points, there exist constants
$C,c>0$, independent of $m$, such that
\[
 \max_{1\le\ell,h\le s}
 \left|
 (\mathbf L_m-\mathbf L_{\infty,m})_{\ell h}
 \right|
 \le C m^4 e^{-c\sqrt m}.
\]
Consequently,
\[
        \frac{1}{m^2}
        \bigl(\mathbf L_m-\mathbf L_{\infty,m}\bigr)
        \longrightarrow0.
\]
\end{lemma}

\begin{proof}
All coefficients in the finite sums defining $D,K,B$ and their
realignments grow at most polynomially, of degree four, in the depth
indices.  Since the set of speeds is fixed and
$\rho_0:=\min_\ell\Re a_\ell>0$, every test-point monomial of total
depth $q$ has modulus at most $e^{-\rho_0 q/\sqrt m}$.  Every term
discarded when the infinite sums are truncated at depth $m$ contains
at least one index of size at least $m$.  Summing the resulting
polynomially weighted geometric tails gives
\[
 C m^4e^{-c\sqrt m}
\]
for suitable constants $C,c>0$.  Division by $m^2$ proves the second
assertion.
\end{proof}

For the fixed speed parameters define
\begin{equation}\label{B:eq:AC-limit}
 A_{\ell h}:=
 |a_\ell+\overline{a_h}|^2,\qquad
 C_{\ell h}:=
 4\,\Re(a_\ell)\Re(a_h).
\end{equation}
Since $1-e^{-\delta z}=\delta z+O(\delta^2)$ as $\delta\to0$,
uniformly for the present finite set of $z$'s,
\[
 m A_m^{\ell h}\longrightarrow A_{\ell h},\qquad
 m C_m^{\ell h}\longrightarrow C_{\ell h},\qquad
 \Delta_m^{\ell h}\longrightarrow4.
\]
Combining this with \eqref{B:eq:Linfty} and
\Cref{B:lem:tail} gives
\[
 \frac{1}{m^2}(\mathbf L_m)_{\ell h}
 \longrightarrow
 2\left(
 \frac{2}{A_{\ell h}^2}
 +\frac{1}{C_{\ell h}^2}
 -\frac{1}{A_{\ell h}C_{\ell h}}
 \right).
\]
Since $s$ is fixed, the convergence is also convergence in every
matrix norm.  This is the limit matrix $\mathbf L_\star$ used in
\Cref{lem:negative-limit}.

\subsection{Qualitative control of the Gram-dependent part}

We now prove the only uniform estimate needed in the main theorem:
for the fixed finite test,
\[
 \sup_{\widehat Q_{\rm av}\succeq0}
 \frac{1}{m^2}
 \max_{\ell,h}|(\mathbf R_m)_{\ell h}|
 \longrightarrow0.
\]
The argument has two ingredients.  Positivity first eliminates all
low-total mixed modes.  This gives an improved estimate on conjugate
points.  A qualitative maximum-principle argument then propagates the
improvement to the finitely many realigned points. 
For $0\le g\le2m-2$, let
\[
 s_g:=\sum_{\substack{p+q=g\\0\le p,q<m}}\mathbf e_{pq}
\]
in the mixed coordinates.

\begin{lemma}[Forced low mixed modes]\label{B:lem:means}
The baseline mixed block $B$ is positive semidefinite, and every
positive semidefinite normal-form Gram representation satisfies
\begin{equation}\label{B:eq:means}
                 Ms_g=Bs_g=0,\qquad 0\le g<m.
\end{equation}
\end{lemma}

\begin{proof}
The matrix $B$ is block diagonal with respect to the total
$g=p+q$.  On a complete block, its off-diagonal entries are
nonpositive and its row sums vanish; hence the block is a weighted
graph Laplacian and is positive semidefinite.  Incomplete blocks are
principal submatrices, so $B\succeq0$.

To obtain the second assertion, substitute the corner variables so that
the pure wedges vanish and the mixed wedge $w_{pq}$ equals $t^{p+q}$.
If $Q\succeq0$ is any feasible Gram matrix, the resulting polynomial
is
\[
 \left\|\sum_g t^gQ^{1/2}s_g\right\|^2.
\]
The same polynomial computed from the baseline mixed block has no term
of degree below $2m$.  If $g<m$ were the least index for which
$Q^{1/2}s_g\ne0$, the coefficient of $t^{2g}$ would be strictly
positive, a contradiction.  Thus $Qs_g=0$ for $g<m$.  Restricting to
the mixed block gives \eqref{B:eq:means}.
\end{proof}

Write
\[
        \Psi_m:=M-B.
\]
By \Cref{B:lem:means}, the repeated generating expression associated
with $\Psi_m$ contains only total degrees at least $m$ in each pair of
variables.  At the conjugate points used below this high-order support
is exponentially suppressed.  Together with the explicit baseline
block $D-K$, this yields the following qualitative estimate.

\begin{lemma}[Conjugate-point improvement]\label{B:lem:conjugate}
Let $\alpha,\beta$ range over fixed compact real intervals with
$\alpha>0$, and set
\[
 \xi_m=ie^{-(\alpha+i\beta)/\sqrt m},\qquad
 \eta_m=-ie^{-(\alpha-i\beta)/\sqrt m}.
\]
Uniformly over all positive semidefinite normal-form Gram
representations,
\[
        P_-[\xi_m,\eta_m]=O(m),
\]
whereas on every fixed complex neighbourhood on which the two node
moduli stay strictly inside the unit disk at distance comparable with
$m^{-1/2}$,
\[
        P_-[\xi,\eta]=O(m^2).
\]
\end{lemma}

\begin{proof}
The global estimate follows from positivity and the diagonal pure
block:
\[
 0\le P_-[\xi,\eta]
 \le
 \frac{8}{(1-|\xi|^2)^2(1-|\eta|^2)^2}.
\]
At the scaled points above, the right-hand side is $O(m^2)$.

On the conjugate slice, the same-side four-form vanishes and the
Pl\"ucker correction can be rewritten through $\Psi_m$.  The baseline
term $D-K$ is $O(m)$ there, while the high-total support of $\Psi_m$
from \Cref{B:lem:means} contributes $o(m)$.  Hence
$P_-[\xi_m,\eta_m]=O(m)$ uniformly.  A direct tail estimate proving the last $o(m)$ statement is obtained exactly as in \Cref{B:lem:tail}; the quantitative version is recorded below.
\end{proof}

We use the following standard qualitative form of the two-constant
principle.

\begin{lemma}[Qualitative continuation]\label{B:lem:continuation}
Let $W_m$ be a holomorphic vector-valued function on a fixed rectangle
$\Omega\subset\C$, possibly with a range dimension depending on $m$.
Assume $\|W_m\|\le C$ on $\Omega$ and
$\|W_m\|\to0$ uniformly on a nontrivial open segment $I$ of one side
of the rectangle.  Then $\|W_m\|\to0$ uniformly on every compact subset
of $\Omega$ having positive harmonic measure with respect to $I$.
\end{lemma}

\begin{proof}
For nonzero $W_m$, the function $\log\|W_m\|$ is subharmonic; zeros are
handled by the usual value $-\infty$.  The two-constant theorem bounds
$\log\|W_m(z)\|$ by the boundary majorant weighted by the harmonic
measure of $I$.  On a fixed compact subset this harmonic measure is
bounded below by a positive constant.  Since the bound on $I$ tends to
zero while the global bound remains fixed, the resulting interior bound
also tends to zero.  The identically zero case is immediate.
\end{proof}

\begin{proposition}[Uniform vanishing of the Gram-dependent term]
\label{B:prop:R}
For the fixed test points of \eqref{eq:test-nodes},
\begin{equation}\label{B:eq:Rvanish}
 \sup_{\widehat Q_{\rm av}\succeq0}
 \frac1{m^2}\max_{1\le\ell,h\le s}
 |(\mathbf R_m)_{\ell h}|
 \longrightarrow0.
\end{equation}
Consequently, for every fixed $c\in\C^s$,
\[
 \sup_{\widehat Q_{\rm av}\succeq0}
 \left|
 c^*\frac{\mathbf R_m}{m^2}c
 \right|
 \longrightarrow0.
\]
\end{proposition}

\begin{proof}
Let $P_-=\Gamma^\top\Gamma$ and define
\[
 \mathbf h_m(\alpha,\beta)
 :=
 \frac1m\,
 \Gamma f\left(
 ie^{-(\alpha+i\beta)/\sqrt m},
 -ie^{-(\alpha-i\beta)/\sqrt m}
 \right).
\]
By \Cref{B:lem:conjugate}, $\|\mathbf h_m\|$ is uniformly bounded on
fixed complex rectangles, while it tends uniformly to zero on the real
conjugate slice.

Apply \Cref{B:lem:continuation} first in the $\alpha$ variable and then
in the $\beta$ variable.  No explicit harmonic-measure constant is
needed because the test set is finite and fixed.  For each pair
$(\ell,h)$, put
\[
 \alpha_{\ell h}:=
 \frac{a_\ell+\overline{a_h}}2,\qquad
 \beta_{\ell h}:=
 \frac{a_\ell-\overline{a_h}}{2i}.
\]
These finitely many points lie in a fixed compact set.  Realignment
gives
\[
 \frac1{m^2}(\mathbf R_m)_{\ell h}
 =
 \left\|
 \mathbf h_m(\alpha_{\ell h},\beta_{\ell h})
 \right\|^2,
\]
and therefore every entry tends to zero, uniformly over all feasible
Gram representations.  Since there are only finitely many entries,
\eqref{B:eq:Rvanish} follows.  Finally,
\[
 \left|
 c^*\frac{\mathbf R_m}{m^2}c
 \right|
 \le
 \|c\|_1^2\,
 \frac1{m^2}\max_{\ell,h}|(\mathbf R_m)_{\ell h}|,
\]
which proves the last assertion.
\end{proof}

\subsection{Explicit estimates for the sufficient threshold}
\label{B:subsec:effective}

For completeness, we now record the quantitative version of the preceding
argument that is used in \Cref{app:explicit-bound}.  These constants are not
needed for the qualitative proof in the main text.

For four arguments in the open unit disk define
\begin{equation}\label{B:eq:denominators}
\begin{gathered}
 A=(1-\xi\zeta)(1-\eta\omega),\qquad
 \Delta=(1-\xi\omega)(1-\eta\zeta),\\
 C=(1-\xi\eta)(1-\zeta\omega),\qquad
 \Pi=\xi\eta\zeta\omega.
\end{gathered}
\end{equation}
At the level of scalar generating expressions, write
$Q_{0,m}:=D-K$ and let $L_m$ denote the Gram-invariant kernel
corresponding to \(\mathbf L_m\).  A subscript \(\infty\) means that the
depth indices are extended over all nonnegative integers.  Direct summation
of the geometric series gives
\begin{equation}\label{B:eq:generating}
\begin{aligned}
 D_\infty&=2(A^{-2}-\Delta^{-2}),\\
 B_\infty&=2(A^{-2}-(A\Delta)^{-1}),\\
 Q_{0,\infty}&=D_\infty-\frac{2(\Delta-A)}{A\Delta C},\\
 L_\infty&=2\left(\frac2{A^2}+\frac1{C^2}
                       -\frac1{AC}-\frac2{\Delta^2}\right).
\end{aligned}
\end{equation}
Indeed, the first two formulas follow from
$\sum_{p\ge0}(p+1)t^p=(1-t)^{-2}$ and the common-shift identity for the
minimum kernel.  For the pure cross block, setting $q=p+h$, $s=r+h$,
$h\ge1$, gives
\[
 \sum_{p,r\ge0}\min(p+1,r+1)(\xi\eta)^p(\zeta\omega)^r
     =\frac1{C(1-\Pi)},
\]
and
\[
 \sum_{h\ge1}(\eta^h-\xi^h)(\omega^h-\zeta^h)
     =\frac{(1-\Pi)(\Delta-A)}{A\Delta}.
\]
Realignment interchanges $A$ and $C$ and leaves $\Delta$ unchanged,
which yields the last line of \eqref{B:eq:generating}.

Write $\mathcal E_{\rm mix}:=M-B$.  By \Cref{B:lem:means}, the repeated
generating expression $\mathcal E_{\rm mix}(\xi,\xi,\zeta,\zeta)$
contains only total degrees at least $m$ in each pair of variables.
Moreover, positivity and \eqref{eq:sectors} imply the global bound
\begin{equation}\label{B:eq:global}
 0\le P_\pm[\xi,\eta]
 \le\frac8{(1-|\xi|^2)^2(1-|\eta|^2)^2},
 \qquad |\xi|,|\eta|<1.
\end{equation}
On conjugate nodes, the same-side four-form vanishes and
\begin{equation}\label{B:eq:repeated}
 P_-[\xi,\bar\xi]
 =Q_{0,m}[\xi,\bar\xi]
  +\mathcal E_{\rm mix}(\xi,\xi,\bar\xi,\bar\xi).
\end{equation}
If $\xi=\sqrt\rho e^{i\vartheta}$, the corresponding infinite baseline is
\begin{equation}\label{B:eq:repeated-baseline}
 Q_{0,\infty}[\xi,\bar\xi]
 =\frac{8\rho\sin^2\vartheta}
 {(1-\rho)^2((1-\rho)^2+4\rho\sin^2\vartheta)^2}.
\end{equation}

\begin{lemma}[Uniform finite tails]\label{B:lem:quant-tails}
For $0<\epsilon\le1$, set
\begin{equation}\label{B:eq:tail-budget}
 \mathcal T(m,\epsilon)=2^{32}\epsilon^{-8}e^{-\epsilon m/4}.
\end{equation}
If all four argument moduli are at most $e^{-\epsilon/2}$, then
$|L_m-L_\infty|\le\mathcal T(m,\epsilon)$.  For every positive
semidefinite normal-form Gram representation,
\begin{equation}\label{B:eq:conjugate-tail}
 |P_-[\xi,\bar\xi]-Q_{0,\infty}[\xi,\bar\xi]|
 \le\mathcal T(m,\epsilon)
 \qquad(|\xi|\le e^{-\epsilon/2}).
\end{equation}
\end{lemma}

\begin{proof}
The diagonal of the realigned correction vanishes, so the positive
semidefinite matrices $M$ and $B$ have the same diagonal.  For
$W=(p+1)(q+1)(r+1)(s+1)$, Cauchy--Schwarz gives
\[
 |(\mathcal E_{\rm mix})_{pq,rs}|\le4\sqrt W\le4W.
\]
The coefficients of the baseline kernels are bounded by a constant
multiple of $W$.  Every discarded finite-depth term contains at least one
index of size at least $m$.  Splitting the radial decay in half and using
$1-e^{-\epsilon/4}\ge\epsilon/8$ therefore gives
\[
 |L_m-L_\infty|,\ |Q_{0,m}-Q_{0,\infty}|
 \le2^{30}\epsilon^{-8}e^{-\epsilon m/4}.
\]
For the correction in \eqref{B:eq:repeated}, \Cref{B:lem:means} leaves
only total degrees at least $m$ in both variable pairs; the same summation
gives
\[
 |\mathcal E_{\rm mix}(\xi,\xi,\bar\xi,\bar\xi)|
 \le2^{26}\epsilon^{-8}e^{-\epsilon m/2}.
\]
Combining the two estimates proves the lemma.
\end{proof}

For the explicit threshold set
\begin{equation}\label{B:eq:dyadic-scale}
 k\ge16,\qquad \epsilon=2^{-k},\qquad m=\epsilon^{-2}=2^{2k}.
\end{equation}
Then
\begin{equation}\label{B:eq:dyadic-tails}
 \mathcal T\le2^{32+8k-2^{k-2}},\qquad
 \mathcal T\le\epsilon^{-2},\qquad
 \epsilon^4\mathcal T\le\epsilon.
\end{equation}

\begin{lemma}[Quantitative continuation]\label{B:lem:quant-continuation}
Let $P_-=\Gamma^\top\Gamma$ be a real Gram factorization, define
$\mathbf h(\xi,\eta)=\Gamma f(\xi,\eta)$, and set
\begin{equation}\label{B:eq:hepsilon}
 \mathbf h_\epsilon(\alpha,\beta)
 =\mathbf h\bigl(ie^{-\epsilon(\alpha+i\beta)},
        -ie^{-\epsilon(\alpha-i\beta)}\bigr).
\end{equation}
For the parameters in \eqref{B:eq:dyadic-scale},
\begin{equation}\label{B:eq:quant-continuation}
 \epsilon^4\norm{\mathbf h_\epsilon(\alpha,\beta)}^2
 \le2^{20}\epsilon^{2^{-17}}
\end{equation}
throughout
\begin{equation}\label{B:eq:speed-region}
 1\le\Re\alpha\le2,\quad |\Im\alpha|\le256,\quad
 |\Re\beta|\le256,\quad |\Im\beta|\le\tfrac12.
\end{equation}
The bound is uniform over the dimension of the Gram factorization.
\end{lemma}

\begin{proof}
From \eqref{B:eq:global} and $1-e^{-\epsilon/2}\ge\epsilon/4$,
\[
 \norm{\mathbf h_\epsilon(\alpha,\beta)}^2
 \le2^{20}\epsilon^{-4}
\]
whenever $\Re\alpha-|\Im\beta|\ge1/4$.  On the real conjugate slice,
\eqref{B:eq:repeated-baseline}, \Cref{B:lem:quant-tails}, and
\eqref{B:eq:dyadic-tails} improve this to
\[
 \norm{\mathbf h_\epsilon(\alpha,\beta)}^2
 \le2^{20}\epsilon^{-2}
\]
for $\alpha\in[1/2,1024+1/2]$ and $\beta\in[-512,512]$.

We apply the two-constant principle twice.  On the upper
$\alpha$-rectangle
$1/2\le\Re\alpha\le1024+1/2$, $0\le\Im\alpha\le512$, use
\[
 h_1(\alpha)=\sin\frac{\pi(\Re\alpha-1/2)}{1024}
 \frac{\sinh(\pi(512-\Im\alpha)/1024)}{\sinh(\pi/2)}.
\]
On $1\le\Re\alpha\le2$ and $0\le\Im\alpha\le256$ one has
$h_1>2^{-12}>2^{-13}$.  Repeating in the lower rectangle yields
\[
 \norm{\mathbf h_\epsilon(\alpha,\beta)}^2
 \le2^{20}\epsilon^{-4+2\cdot2^{-13}}
\]
for $1\le\Re\alpha\le2$, $|\Im\alpha|\le256$, and real
$\beta\in[-512,512]$.

For fixed such $\alpha$, apply the same argument in the upper and lower
$\beta$-rectangles with
\[
 h_2(\beta)=\cos\frac{\pi\Re\beta}{1024}
 \frac{\sinh(\pi(3/4-\Im\beta)/1024)}{\sinh(3\pi/4096)}.
\]
On $|\Re\beta|\le256$ and $|\Im\beta|\le1/2$, its relevant harmonic
measure is larger than $2^{-5}$.  Since
$2\cdot2^{-13}\cdot2^{-5}=2^{-17}$, the second comparison gives
\eqref{B:eq:quant-continuation}.
\end{proof}


\section{An exact negative certificate for $\mathbf L_\star$}
\label{app:negative-witness}

This appendix certifies the only finite computation used in
\Cref{lem:negative-limit}: the fixed matrix $\mathbf L_\star$ has a
strictly negative eigenvalue.  The construction uses exactly the two
vertical families of test speeds listed in the main text.

For the certificate only, it is convenient to parameterize the same
finite set by
\[
 I=\{(r,j,\sigma):r\in\{1,2\},\ 1\le j\le256,\
 \sigma\in\{-1,1\}\},
\]
with
\[
 a_{r,j,\sigma}=r+i\sigma j.
\]
For $\ell=(r,j,\sigma)$ and $h=(r',j',\sigma')$, the entries of
$\mathbf L_\star$ are
\begin{equation}\label{C:eq:entry}
 (\mathbf L_\star)_{\ell h}
 =
 2\left(
 \frac{2}{A_{\ell h}^2}
 +\frac{1}{C_{\ell h}^2}
 -\frac{1}{A_{\ell h}C_{\ell h}}
 \right),
\end{equation}
where
\[
 A_{\ell h}
 =(r+r')^2+(\sigma j-\sigma'j')^2,\qquad
 C_{\ell h}=4rr'.
\]
Thus every entry is rational.

Define the integer vector $d\in\mathbb Z^{1024}$ by
\begin{equation}\label{C:eq:witness}
 d_{r,j,\sigma}
 :=
 \sigma\widehat d_j,\qquad
 \widehat d_j
 :=
 \left\lfloor
 \frac{2^{43}j}{(128^2+j^2)^2}
 \right\rfloor.
\end{equation}

\begin{lemma}[Exact negative direction]\label{C:lem:witness}
The vector in \eqref{C:eq:witness} satisfies
\begin{equation}\label{C:eq:witness-bounds}
 \|d\|_1=859652168<2^{30},
 \qquad
 d^\top\mathbf L_\star d<-2^{42}.
\end{equation}
Consequently,
\begin{equation}\label{C:eq:eig-margin}
 \lambda_{\min}(\mathbf L_\star)<-2^{-18}.
\end{equation}
In particular, the constant $\gamma$ in
\Cref{lem:negative-limit} may be chosen as $\gamma=2^{-18}$.
\end{lemma}

\begin{proof}
Since the entries in \eqref{C:eq:entry} are rational, the quadratic
form can be checked using integer arithmetic only.  Define
\[
 \mathcal J
 :=
 \sum_{\ell,h\in I}
 \left\lfloor
 \frac{2^{64}d_\ell d_h
 \bigl(4C_{\ell h}^2+2A_{\ell h}^2
       -2A_{\ell h}C_{\ell h}\bigr)}
 {A_{\ell h}^2C_{\ell h}^2}
 \right\rfloor.
\]
Each summand is rounded toward $-\infty$.  Since $|I|^2=2^{20}$,
\[
 \frac{\mathcal J}{2^{64}}
 \le d^\top\mathbf L_\star d
 \le \frac{\mathcal J+2^{20}}{2^{64}}.
\]
The exact integer evaluation is
\[
\begin{aligned}
 \mathcal J
 &=-111807741093765987869670624504700,\\
 \mathcal J+2^{20}
 &=-111807741093765987869670623456124,
\end{aligned}
\]
whose upper numerator is smaller than $-2^{106}$.  Hence
$d^\top\mathbf L_\star d<-2^{42}$.  The stated $\ell_1$ norm follows
directly from \eqref{C:eq:witness}.

Now
\[
 \|d\|_2^2\le\|d\|_1^2<2^{60}.
\]
Therefore the Rayleigh quotient satisfies
\[
 \frac{d^\top\mathbf L_\star d}{\|d\|_2^2}
 <-\frac{2^{42}}{2^{60}}
 =-2^{-18}.
\]
By the variational characterization of the smallest eigenvalue,
\[
 \lambda_{\min}(\mathbf L_\star)
 \le
 \frac{d^\top\mathbf L_\star d}{\|d\|_2^2}
 <-2^{-18},
\]
which proves \eqref{C:eq:eig-margin}.
\end{proof}

\begin{remark}
A direct floating-point diagonalization gives
$
\lambda_{\min}(\mathbf L_\star)\approx-1.18428\times10^{-2},
$
which is substantially more negative than the conservative exact
margin in \eqref{C:eq:eig-margin}.  The numerical value is not used in
the proof.
\end{remark}

\begin{remark}[Alternative families of test points]
\label{C:rem:alternative-test-families}
The two-family choice $\Re a=1,2$ is not unique.  Preliminary numerical
experiments with three vertical families
\[
 \Re a\in\{1,2,3\},\qquad
 \Im a=\pm1,\ldots,\pm J,
\]
show a negative eigenvalue already near $J=44$, corresponding to
$s=6J=264$ test points.  We retain the two-family construction because
it is the family for which the exact certificate above is supplied.
\end{remark}

\section{An explicit finite threshold}\label{app:explicit-bound}

The main proof needs only the qualitative limits in \Cref{B:prop:R}. Here we retain the explicit constants from the quantitative argument and obtain a sufficient finite order.

\begin{theorem}[Uniform finite-scale non-SoS test]\label{thm:finite-scale}
Let $k,\epsilon,m$ satisfy \eqref{B:eq:dyadic-scale}. If $\calF_m$ is SoS,
then every positive semidefinite normal-form Gram representation \eqref{eq:general-gram-1} gives a
nonnegative scalar $\Phi_m$ satisfying
\begin{equation}\label{eq:finite-scale}
          0\le\Phi_m
          <-2^{42}+2^{90}m^{-1/2}+2^{80}m^{-1/262144}.
\end{equation}
All constants are independent of the Gram corrections and of the rank
of the representation.
\end{theorem}
\begin{proof}

Use the fixed data of \eqref{C:eq:witness}, and put
$\mathsf x_\ell=(\xi_\ell,\eta_\ell)
=(ie^{-\epsilon a_\ell},-ie^{-\epsilon\bar a_\ell})$.
Then $\eta_\ell=\bar\xi_\ell$.
Define
\begin{equation}\label{eq:evaluated-gram}
 \mathcal G_m=
 \left[
 \left(M+\frac{P_++P_-}{2}\right)(\overline{\mathsf x_\ell},\mathsf x_h)
 \right]_{\ell,h\in I},
 \qquad \Phi_m=\epsilon^4d^*\mathcal G_md.
\end{equation}
Each summand is an evaluated positive semidefinite kernel on the same
list of points. Taking the direct sum of the pure and mixed feature
spaces proves $\mathcal G_m\succeq0$, hence $\Phi_m\ge0$.

Let $\mathbf L_m$ and $\mathbf R_m$ be the evaluation matrices of
$L_m$ and $\calR P_-$ at these points. Realignment gives
\begin{equation}\label{eq:realigned-diagonal}
 (\calR P_-)(\overline{\mathsf x_\ell},\mathsf x_h)
 =P_-(\bar\xi_\ell,\xi_h,\xi_\ell,\bar\xi_h)
 =P_-[\bar\xi_\ell,\xi_h]=P_-[\xi_\ell,\bar\xi_h].
\end{equation}
The last equality uses real symmetry of $P_-$. For $\ell=(r,j,\sigma)$ and $h=(r',j',\sigma')$, take
\[
 \alpha=\frac{r+r'}2+i\frac{\sigma j-\sigma'j'}2,
 \qquad
 \beta=\frac{\sigma j+\sigma'j'}2-i\frac{r-r'}2.
\]
Then $\alpha+i\beta=a_\ell$ and $\alpha-i\beta=\bar a_h$.
These parameters satisfy \eqref{B:eq:speed-region}, so
\Cref{B:lem:quant-continuation} applies to every entry of
\eqref{eq:realigned-diagonal}. Consequently,
\begin{equation}\label{eq:realigned-error}
 |\epsilon^4d^*\mathbf R_md|
 \le\norm d_1^2\,2^{20}\epsilon^{2^{-17}}
 <2^{80}\epsilon^{2^{-17}}.
\end{equation}

It remains to compare the finite baseline with $\mathbf L_\star$. At the test
points, each factor of $A$ and $C$ has the form $1-e^{-\epsilon\chi}$,
where $\Re\chi\ge2$ and $|\chi|<2^{10}$.
The integral identity
\[
 \frac{1-e^{-\epsilon\chi}}{\epsilon\chi}
 =\int_0^1e^{-\epsilon\chi u}\,du
\]
and $|e^{-\epsilon\chi u}-1|\le\epsilon|\chi|u$ give a relative
error at most $2^{10}\epsilon$ for each factor. Thus
$A/\epsilon^2=A_0(1+\delta_A)$ and
$C/\epsilon^2=C_0(1+\delta_C)$, where
$|\delta_A|,|\delta_C|\le2^{12}\epsilon\le1/16$.

For $|\delta|\le1/4$ we use
$|(1+\delta)^{-1}-1|\le2|\delta|$ and
$|(1+\delta)^{-2}-1|\le6|\delta|$.
At our points $A_0,C_0\ge4$. Each factor of $\Delta$ is
$1+e^{-\epsilon\chi}$, and $|\Im(\epsilon\chi)|\le512\epsilon<1$,
so the exponential has positive real part and $|\Delta|\ge1$.
Writing $\delta=\max(|\delta_A|,|\delta_C|)$, the product of the two
inverse factors has relative error at most $5\delta$. The three reciprocal
terms in \eqref{B:eq:generating} therefore contribute at most
$(3/2+3/4+5/8)\delta<3\delta\le3\cdot2^{12}\epsilon$; the remaining
$\Delta$ term contributes at most $4\epsilon^4$. In particular,
\[
 |\epsilon^4 L_\infty(\overline{\mathsf x_\ell},\mathsf x_h)-\mathbf L_\star(\ell,h)|
 \le2^{20}\epsilon.
\]
All four node moduli are at most $e^{-\epsilon}$, so
\Cref{B:lem:quant-tails} and \eqref{B:eq:dyadic-tails} add at most
$\epsilon$ to this scaled error. We obtain
\begin{equation}\label{eq:finite-tangent}
 |\epsilon^4 L_m(\overline{\mathsf x_\ell},\mathsf x_h)-\mathbf L_\star(\ell,h)|
 \le2^{30}\epsilon.
\end{equation}

Pairing \eqref{eq:finite-decomposition} with $d$, and using
\eqref{C:eq:witness-bounds}, \eqref{eq:realigned-error} and
\eqref{eq:finite-tangent}, gives
\begin{align*}
 0\le\Phi_m
 &=\epsilon^4d^*\mathbf L_md-\epsilon^4d^*\mathbf R_md\\
 &<-2^{42}+2^{90}\epsilon+2^{80}\epsilon^{2^{-17}}\\
 &=-2^{42}+2^{90}m^{-1/2}+2^{80}m^{-1/262144}.
\end{align*}
Here both entrywise errors are multiplied by
$\norm d_1^2<2^{60}$, and $\epsilon=m^{-1/2}$.
This proves \eqref{eq:finite-scale}.
\end{proof}

\begin{proposition}[Explicit sufficient order]\label{prop:explicit-threshold}
Let
\[
 k=38\cdot2^{17}+1=\ScaleExponent,\qquad
 m=2^{2k}=\CornerDepth.
\]
Then $\mathcal F_m$ is not SoS. Consequently, $F_N$ is not SoS for
every
\[
                       N\ge2m=\SufficientOrder.
\]
\end{proposition}
\begin{proof}
Convexity of $x\mapsto2^{-x}$ gives
$2^{-x}\le1-x/2$ for $0\le x\le1$. The upper bound in
\eqref{eq:finite-scale} is therefore at most
\[
 -2^{42}+2^{90-k}+2^{42}2^{-2^{-17}}
 \le-2^{24}+2^{90-k}<0.
\]
This contradicts $\Phi_m\ge0$, so $\calF_m$ is not SoS. By
\Cref{lem:corner}, it is a linear restriction of $F_N$ for every
$N\ge2m=\SufficientOrder$, and any SoS representation of such an
$F_N$ would restrict to one of $\calF_m$.
\end{proof}

The size of this threshold comes from the slow term
$2^{80}m^{-1/262144}$ in \eqref{eq:finite-scale}. The witness itself is
fixed; the large order is needed only to make the uniform continuation
error smaller than its negative margin. This bound is sufficient and is
not intended to identify the first non-SoS order.

\section{Explicit certificates for the positive subclasses}
\label{app:positive-certificates}

This appendix supplies the reduction and the coefficient formulas used
in Proposition~\ref{prop:same-symmetry}. We first compute the
commutator, then rearrange its window terms by Pl\"ucker identities.
All index sets are finite; a sum over an empty set is zero.

\subsection{Reduction for pairs with the same symmetry}
\label{app:same-symmetry-reduction}

For symmetric Toeplitz matrices write $X_{ij}=x_{|i-j|}$ and
$Y_{ij}=y_{|i-j|}$. For skew-symmetric Toeplitz matrices write
$X_{ij}=\operatorname{sgn}(i-j)x_{|i-j|}$ and similarly for $Y$,
with $x_0=y_0=0$. In both cases $z_{ab}=x_ay_b-x_by_a$.
In the symmetric case the diagonal multiplicities are $w_0=N$ and
$w_a=2(N-a)$ for $a>0$. The weighted Lagrange identity
\cite[Lemma~10]{laszlo2012} gives
\[
 2\norm X_F^2\norm Y_F^2-2\langle X,Y\rangle_F^2
 =4N\sum_{b=1}^{N-1}(N-b)z_{0b}^2
   +8\sum_{1\le a<b<N}(N-a)(N-b)z_{ab}^2.
\]
The same formula holds in the skew-symmetric case with its first sum
equal to zero.

Use matrix indices $0,\ldots,N-1$, set $m=N-1$, and fix
$\delta=j-i>0$. For a symmetric pair, splitting the summation index
in $[X,Y]_{ij}$ at $i$ and $j$ gives
\[
 [X,Y]_{i,i+\delta}
 =\sum_{t=0}^{i}z_{t,t+\delta}
  -\sum_{t=0}^{m-\delta-i}z_{t,t+\delta}.
\]
The terms with an index strictly between $i$ and $j$ cancel in pairs.
For a skew-symmetric pair the displayed expression changes sign.
Both commutators are skew-symmetric. Each nonempty window occurs
twice above the diagonal, so
\[
 \norm{[X,Y]}_F^2
 =4\sum_{\delta=1}^{m-1}
   \sum_{\substack{k\ge0\\2k<m-\delta}}
   \left(\sum_{t=k+1}^{m-\delta-k}z_{t,t+\delta}\right)^2.
\]
Reflection $t=N-\delta-p$ turns each window into the corresponding
sum of $\hat z_{p,p+\delta}=z_{N-p-\delta,N-p}$. Combining these
identities proves \eqref{eq:positive-reduction}.

\subsection{The window certificate}\label{app:positive-windows}

For $1\le\delta<m$, put $h_\delta=m-\delta$ and define
\[
 \kappa_\delta(p,q)
 =\min(p,q,h_\delta+1-p,h_\delta+1-q),
 \qquad 1\le p,q\le h_\delta.
\]
This counts the nested intervals
$[k+1,h_\delta-k]$ containing both $p$ and $q$. Consequently,
\begin{equation}\label{eq:positive-tent}
 \sum_{\substack{k\ge0\\2k<h_\delta}}
 \left(\sum_{p=k+1}^{h_\delta-k}v_p\right)^2
 =\sum_{p,q=1}^{h_\delta}\kappa_\delta(p,q)v_pv_q.
\end{equation}
Thus the uncorrected Gram of $\Psi_m$ has diagonal
$2p(p+\delta)-\kappa_\delta(p,p)$ and entries
$-\kappa_\delta(p,q)$ between distinct wedges of gap $\delta$.

For $3\le\sigma\le2m-1$, set
\[
 \begin{gathered}
 I_\sigma=\{a\in\mathbb Z:\max(1,\sigma-m)\le a<\sigma/2\},\\
 c_\sigma=\max(0,\sigma-m-1),\qquad
 L_\sigma=\min(\sigma,m+1),\\
 \mathcal W_k^\sigma
 =\{a\in I_\sigma:c_\sigma+1+k\le a\le L_\sigma-1-k\},\\
 t_\sigma(a,b)
 =\min(a-c_\sigma,b-c_\sigma,L_\sigma-a,L_\sigma-b).
 \end{gathered}
\]
For $1\le p<q\le h_\delta$ with $e=q-p\le\delta$, define
\[
 \omega_\delta(p,q)=
 \begin{cases}
 \kappa_\delta(p,q)+\kappa_e(p,p+\delta),&e<\delta,\\
 \kappa_\delta(p,q),&e=\delta.
 \end{cases}
\]
The remaining diagonal coefficient is
\begin{equation}\label{eq:positive-slack}
 s_\delta(p)=2p(p+\delta)-\kappa_\delta(p,p)
 -t_{2p+\delta}(p,p)
 -\sum_{\substack{1\le q\le h_\delta\\0<|q-p|\le\delta}}
 \omega_\delta(\min(p,q),\max(p,q)).
\end{equation}

\begin{proposition}[An explicit window decomposition]
\label{prop:positive-window-certificate}
For every $m\ge1$, the reflected wedges in
\eqref{eq:positive-reduced-form} satisfy
\begin{equation}\label{eq:positive-window-certificate}
 \begin{aligned}
 \Psi_m(\hat z)
 &=\sum_{\sigma=3}^{2m-1}\sum_{k\ge0}
     \left(\sum_{a\in\mathcal W_k^\sigma}
                  \hat z_{a,\sigma-a}\right)^2 +\sum_{\delta=1}^{m-1}
   \sum_{\substack{1\le p<q\le h_\delta\\q-p\le\delta}}
   \omega_\delta(p,q)
       (\hat z_{p,p+\delta}-\hat z_{q,q+\delta})^2 +\sum_{\delta=1}^{m-1}\sum_{p=1}^{h_\delta}
      s_\delta(p)\hat z_{p,p+\delta}^2.
 \end{aligned}
\end{equation}
All nonempty edge weights are positive integers, and
$s_\delta(p)\ge p^2$. In particular, the right-hand side is SoS
in the original symbol variables.
\end{proposition}

\begin{proof}

For $p\ge1$, $1\le\alpha<\beta$ and
$p+\alpha+\beta\le m$, the Pl\"ucker identity is
\[
 R_{p,\alpha,\beta}
 :=\hat z_{p,p+\alpha}\hat z_{p+\beta,p+\alpha+\beta}
  -\hat z_{p,p+\beta}\hat z_{p+\alpha,p+\alpha+\beta}
  +\hat z_{p,p+\alpha+\beta}\hat z_{p+\alpha,p+\beta}=0.
\]
Add $2\kappa_\alpha(p,p+\beta)R_{p,\alpha,\beta}$ to
$\Psi_m$ for every such triple. These additions preserve the
polynomial. Their effect on the Gram entries is as follows.
For two wedges of the same gap $\delta$ at positions $p<q$,
write $e=q-p$. If $e>\delta$, the correction cancels the entry
$-\kappa_\delta(p,q)$. If $e<\delta$, it changes that entry to
$-\kappa_\delta(p,q)-\kappa_e(p,p+\delta)=-\omega_\delta(p,q)$.
If $e=\delta$, no four distinct indices are present and the entry
remains $-\omega_\delta(p,q)$.

The third product in $R_{p,\alpha,\beta}$ couples wedges on the
same antidiagonal. For $a<b$ in $I_\sigma$, its coefficient in
the Gram is
\[
 \kappa_{b-a}(a,\sigma-b)
 =\min(a,\sigma-b,m+1-b,m+1+a-\sigma)
 =t_\sigma(a,b).
\]
Counting the sets $\mathcal W_k^\sigma$ containing both $a$ and
$b$ expresses these entries by the first sum in
\eqref{eq:positive-window-certificate}. The second sum gives
the negative same-gap entries, and \eqref{eq:positive-slack}
makes the diagonals agree. All remaining entries are zero.
This proves the identity.

It remains to check the signs. Each tent value in its range is
a positive integer, so $\omega_\delta(p,q)\ge1$.
The bounds $\kappa_\delta(p,p)\le p$ and
$t_{2p+\delta}(p,p)\le p$ are immediate.
For a right neighbour at distance $e<\delta$ the edge weight
is at most $2p$, and at distance $e=\delta$ it is at most $p$.
For a left neighbour these bounds are $2(p-e)$ and $p-\delta$,
respectively. Including missing boundary neighbours only increases
the upper bound, whence
\[
 \sum_q\omega_\delta(\min(p,q),\max(p,q))
 \le
 \begin{cases}
 (2\delta-1)p+p(p-1),&p\le\delta,\\
 (4\delta-2)p-\delta^2,&p>\delta.
 \end{cases}
\]
Inserting these estimates into \eqref{eq:positive-slack} gives
\[
 s_\delta(p)\ge
 \begin{cases}
 p^2,&p\le\delta,\\
 p^2+(p-\delta)^2,&p>\delta.
 \end{cases}
\]
The sums in the certificate are finite because the windows are
eventually empty. Every wedge is quadratic in the original variables,
so the displayed identity is a finite sum of homogeneous quadratic
squares. Integer weights may equivalently be realized by repeating
the corresponding squares.
\end{proof}

\subsection{The coupling when both factors have both parts}
\label{app:positive-coupling}

Let $X=S_1+K_1$ and $Y=S_2+K_2$, with $S_1,S_2$ symmetric and $K_1,K_2$
skew-symmetric. Expanding the two Frobenius-orthogonal parts of
the commutator gives
\[
 F(X,Y)=F(S_1,S_2)+F(K_1,K_2)+F(S_1,K_2)+F(K_1,S_2)+\mathcal C,
\]
where
\[
 \begin{aligned}
 \mathcal C= -4\langle S_1,S_2\rangle_F\langle K_1,K_2\rangle_F
 -2\langle[S_1,S_2],[K_1,K_2]\rangle_F -2\langle[S_1,K_2],[K_1,S_2]\rangle_F.
 \end{aligned}
\]
For Toeplitz matrices the four displayed $F$ terms have the
certificates of Section~\ref{sec:positive-subclasses}. The term
$\mathcal C$ vanishes if one factor has only one symmetry part,
but is present in general. No sign or SoS assertion for
$\mathcal C$ is needed in the one-sided proof.

\section{Data and formal verification}
\label{app:finite-certificates}

This appendix specifies the $48$ exact positive-increment steps used in
Lemma~\ref{lem:finite-increments} and the exact tests that establish
their validity. The accompanying files
\path{verification/data/step_02.json.gz} through
\path{verification/data/step_49.json.gz} contain rational
coefficients and their indices. The manifest lists the steps in
order and gives hashes of both compressed and uncompressed data.

\subsection{The two types of increment}\label{app:increment-data}

For $2\le n\le16$, the certificate represents
\begin{equation}\label{eq:appendix-normalized-step}
 R_n=(n-1)F_{n+1}-(n+1)F_n
\end{equation}
as an SoS. In \eqref{eq:finite-chain} this gives
$\alpha_n=(n+1)/(n-1)$, the coordinate projection for $\psi_n$,
and the Gram of $R_n$ divided by $n-1$ for $\mathbf H_n$.
For $17\le n\le49$, it represents
\begin{equation}\label{eq:appendix-geometric-step}
 I_n=F_{n+1}-F_n\circ\psi_n,\qquad
 \psi_n(x,y)=
 \bigl((q_n^{|a|}x_a)_{|a|<n},(q_n^{|a|}y_a)_{|a|<n}\bigr),
\end{equation}
where $q_n=1-\theta_n/n$ is rational and $0<q_n\le1$.
The rational number $\theta_n$ is stored in the corresponding data
file. Here $\alpha_n=1$. In particular, $q_{49}=97/98$ at the final step.

The stored Gram blocks use noncentral wedges, with both labels
nonzero. Central wedges are added as separate squares. For each
$1\le a\le n$, the coefficient of each of
$z_{-a,0}^2$ and $z_{0,a}^2$ is
\[
 \begin{cases}
 2(n+1)(a-1),&\text{in }R_n,\\
 2\bigl((n+1)(n+1-a)-q_n^{2a}n(n-a)\bigr),
     &\text{in }I_n.
 \end{cases}
\]
These coefficients are nonnegative. Adjoining them produces a
Gram on all wedges, as required in \eqref{eq:finite-chain}.

\subsection{Exact polynomial and positivity checks}
\label{app:exact-psd-checks}

For each step, the checker expands $F_n$ and $F_{n+1}$ directly
from the Toeplitz matrix definition. It also expands the proposed
Gram expression and verifies equality of every monomial coefficient.
The Pl\"ucker corrections are separately checked to represent zero.
All coefficients in these comparisons are integers or rationals.

The full noncentral Gram has verified kernel vectors
\[
 h_g=\sum_{\substack{-n\le a<0<b\le n\\b-a=g}}e_{ab},
 \qquad n+1\le g\le2n,
\]
where $e_{ab}$ is the coordinate vector for $z_{ab}$.
Their supports are disjoint. For each $h_g$, delete the first
coordinate in its support, in lexicographic order. The remaining
coordinate vectors together with the $h_g$ form a basis. Once
the exact relations $\mathbf H_nh_g=0$ have been checked, this change of
basis gives a congruence
\[
 \mathbf H_n\ \sim\ \begin{pmatrix}\widehat H_n&0\\0&0\end{pmatrix},
\]
where $\widehat H_n$ is the retained principal submatrix.
Consequently $\mathbf H_n\succeq0$ follows from $\widehat H_n\succeq0$.
No theorem about the kernel of all possible Grams is needed for
this check of the supplied matrices.

Each connected block of $\widehat H_n$ is checked independently.
One available test is Sylvester's criterion using exact
fraction-free elimination. The default checker uses an integer
residual certificate. Write a rational block as $A/d$, with
$A$ integral and $d>0$. For an integer matrix $L$ and integer
$s\ge0$, it computes
\begin{equation}\label{eq:exact-residual-certificate}
 E=2^{2s}A-dLL^\top,\qquad
 \delta=\min_i\left(E_{ii}-\sum_{j\ne i}|E_{ij}|\right).
\end{equation}
If $\delta>0$, symmetry and diagonal dominance give
$E\succeq\delta I$, hence
\[
 \frac{A}{d}
 =2^{-2s}LL^\top+\frac{E}{d\,2^{2s}}
 \succeq\frac{\delta}{d\,2^{2s}}I\succ0.
\]
The proposed $L$ may be obtained numerically, but every equality
and inequality in this acceptance test uses exact integers.
The rational factorization coefficients of an SoS itself are
not required: a positive semidefinite rational Gram admits the
real factorization used in Section~\ref{sec:exact-certificates}.

Run \texttt{python3 verification/positive.py} after installing
the dependencies in \path{verification/requirements.txt}.
The checker validates the complete chain, including the base case,
coefficient identities, central squares, kernel relations, and
positivity. It also rejects deliberately corrupted identities
and Gram corrections. The expected final line is
\begin{quote}\small\ttfamily
PASS: F\_N is SOS for 2 <= N <= 50; 48 exact steps; negative controls rejected
\end{quote}
The compressed coefficients are included with the source archive;
the mathematical acceptance conditions are
\eqref{eq:appendix-normalized-step}--\eqref{eq:exact-residual-certificate}.

\subsection{Lean verification of the large-order negative result}
\label{app:lean-negative}

\rev{The public repository also contains a Lean~4 formalization of the
large-order negative theorem for the original Toeplitz BW polynomial.
In particular, the declaration
\texttt{ToeplitzSOS.Negative.sharpNegativeResolution} proves that the
original real Toeplitz BW form is not SoS for every
$N\ge\SufficientOrder$, where the SoS predicate allows an arbitrary
finite number of homogeneous quadratic squares with arbitrary real
coefficients.  The source is
\path{lean/ToeplitzSOS/Negative/Resolution.lean}.  The verified
formalization is frozen at repository tag \texttt{v1.0}; reproduction
commands and the exact verified source tree are recorded in
\path{VERIFICATION.md}.}

\subsection{The positive Lean statement}\label{app:positive-lean}

In the supplied formalization, \texttt{toeplitzBW n} is the
polynomial defined by the original real Toeplitz matrices, and
\texttt{IsSumSqHomQuad} asserts the existence of a finite sum of
real homogeneous quadratic squares. In the namespace
\texttt{ToeplitzSOS}, the range theorem is
\begin{lstlisting}[basicstyle=\ttfamily\small,breaklines=true]
theorem toeplitzBW_isSumSq_of_le_20
    (n : Nat) (h2 : 2 <= n) (h20 : n <= 20) :
    IsSumSqHomQuad (toeplitzBW n)
\end{lstlisting}
Its proof combines the explicit order-two identity with the
fixed-order theorems for $n=3,\ldots,20$.
The source file in the public repository is
\path{lean/ToeplitzSOS/Certificates/Range.lean}.
The Lean and Mathlib versions are pinned by the repository release.
From the \path{lean/} directory, the positive certificates can be
checked with \texttt{./verify --certificates}.  The rational chain at
orders $21$--$50$ is not part of this theorem.


\scriptsize
\bibliographystyle{plainnat}
\bibliography{references}
\end{document}
